\documentclass[11pt]{article}
\usepackage[margin=1.1in]{geometry}
\usepackage{amsmath,amssymb,amsthm,mathtools,booktabs,array,tabularx,longtable}
\usepackage[dvipsnames]{xcolor}
\usepackage[colorlinks=true,linkcolor=MidnightBlue,citecolor=MidnightBlue,urlcolor=MidnightBlue]{hyperref}
\usepackage{fancyhdr}
\pagestyle{fancy}\fancyhf{}\renewcommand{\headrulewidth}{0pt}
\fancyfoot[R]{\thepage}
\fancyfoot[L]{\scriptsize\color{gray}Generated by the Physics Graph, not the version of record. Official version: \href{https://doi.org/10.5281/zenodo.24403592}{doi:10.5281/zenodo.24403592}}
\fancypagestyle{plain}{\fancyhf{}\fancyfoot[R]{\thepage}\fancyfoot[L]{\scriptsize\color{gray}Generated by the Physics Graph, not the version of record. Official version: \href{https://doi.org/10.5281/zenodo.24403592}{doi:10.5281/zenodo.24403592}}\renewcommand{\headrulewidth}{0pt}}
\newtheorem{theorem}{Theorem}[section]
\newtheorem{proposition}[theorem]{Proposition}
\newtheorem{lemma}[theorem]{Lemma}
\newtheorem{corollary}[theorem]{Corollary}
\newtheorem{conjecture}[theorem]{Conjecture}
\newtheorem{refutation}[theorem]{Refutation}
\newtheorem{prediction}[theorem]{Prediction}
\newtheorem{principle}[theorem]{Principle}
\newtheorem{observation}[theorem]{Observation}
\newtheorem{axiom}[theorem]{Axiom}
\theoremstyle{definition}
\newtheorem{definition}[theorem]{Definition}
\newtheorem{fact}[theorem]{Fact}
\newtheorem{assembled}[theorem]{Assembly}
\newtheorem{postulate}[theorem]{Postulate}
\newtheorem{premise}[theorem]{Premise}
\newtheorem{hypothesis}[theorem]{Hypothesis}
\newtheorem{remark}[theorem]{Remark}
\newtheorem{example}[theorem]{Example}
\theoremstyle{plain}
\setlength{\marginparwidth}{0.8in}
\newcommand{\changedmark}[2]{\marginpar{\raggedright\scriptsize\color{#1}#2}}
\providecommand{\Aa}{}\renewcommand{\Aa}{\mathcal{A}}
\providecommand{\Ad}{}\renewcommand{\Ad}{\widetilde{\mathrm{Ad}}}
\providecommand{\Aev}{}\renewcommand{\Aev}{\mathcal{A}^{+}}
\providecommand{\Aod}{}\renewcommand{\Aod}{\mathcal{A}^{-}}
\providecommand{\B}{}\renewcommand{\B}{\mathrm{B}}
\providecommand{\C}{}\renewcommand{\C}{\mathbb{C}}
\providecommand{\CAR}{}\renewcommand{\CAR}{\mathrm{CAR}}
\providecommand{\CS}{}\renewcommand{\CS}{{\mathbb{S}^{+}}}
\providecommand{\CSm}{}\renewcommand{\CSm}{{\mathbb{S}^{-}}}
\providecommand{\Cl}{}\renewcommand{\Cl}{\mathrm{Cl}}
\providecommand{\Con}{}\renewcommand{\Con}{\mathcal{C}}
\providecommand{\Dl}{}\renewcommand{\Dl}{\Delta}
\providecommand{\EN}{}\renewcommand{\EN}{E_{N}}
\providecommand{\End}{}\renewcommand{\End}{\mathrm{End}}
\providecommand{\F}{}\renewcommand{\F}{\mathbb{F}}
\providecommand{\G}{}\renewcommand{\G}{\Gamma}
\providecommand{\HH}{}\renewcommand{\HH}{\mathbb{H}}
\providecommand{\IS}{}\renewcommand{\IS}{I\sigma_3}
\providecommand{\J}{}\renewcommand{\J}{\mathbf{J}}
\providecommand{\LH}{}\renewcommand{\LH}{L_{H}}
\providecommand{\Mat}{}\renewcommand{\Mat}{\mathrm{M}}
\providecommand{\Mo}{}\renewcommand{\Mo}{M^{\circ}}
\providecommand{\Mp}{}\renewcommand{\Mp}{M^{+}}
\providecommand{\Msun}{}\renewcommand{\Msun}{M_{\odot}}
\providecommand{\Mt}{}\renewcommand{\Mt}{\widetilde{M}}
\providecommand{\Ncone}{}\renewcommand{\Ncone}{\mathcal{N}}
\providecommand{\Neff}{}\renewcommand{\Neff}{N_{\mathrm{eff}}}
\providecommand{\PCS}{}\renewcommand{\PCS}{\mathbb{P}{\mathbb{S}^{+}}}
\providecommand{\PT}{}\renewcommand{\PT}{\mathbb{PT}}
\providecommand{\Pin}{}\renewcommand{\Pin}{\mathrm{Pin}}
\providecommand{\Q}{}\renewcommand{\Q}{\mathbb{Q}}
\providecommand{\R}{}\renewcommand{\R}{\mathbb{R}}
\providecommand{\Ric}{}\renewcommand{\Ric}{\mathrm{Ric}}
\providecommand{\SBH}{}\renewcommand{\SBH}{S_{\mathrm{BH}}}
\providecommand{\Sc}{}\renewcommand{\Sc}{\mathbb{S}}
\providecommand{\Sch}{}\renewcommand{\Sch}{\mathrm{P}}
\providecommand{\Sig}{}\renewcommand{\Sig}{\Sigma}
\providecommand{\Som}{}\renewcommand{\Som}{S_\omega}
\providecommand{\Sp}{}\renewcommand{\Sp}{\mathbb{S}}
\providecommand{\Sph}{}\renewcommand{\Sph}{\mathrm{S}}
\providecommand{\Spi}{}\renewcommand{\Spi}{S_\pi}
\providecommand{\Spin}{}\renewcommand{\Spin}{\mathrm{Spin}}
\providecommand{\Sth}{}\renewcommand{\Sth}{S_{\mathrm{th}}}
\providecommand{\T}{}\renewcommand{\T}{\mathbb{T}}
\providecommand{\TH}{}\renewcommand{\TH}{T_{\mathrm{H}}}
\providecommand{\Tr}{}\renewcommand{\Tr}{\operatorname{Tr}}
\providecommand{\UHF}{}\renewcommand{\UHF}{\mathrm{UHF}}
\providecommand{\Z}{}\renewcommand{\Z}{\mathbb{Z}}
\providecommand{\Zc}{}\renewcommand{\Zc}{\mathrm{Z}}
\providecommand{\cmg}{}\renewcommand{\cmg}{\mathrm{cm^{2}\,g^{-1}}}
\providecommand{\dS}{}\renewcommand{\dS}{\mathrm{dS}}
\providecommand{\dsc}{}\renewcommand{\dsc}{\mathrm{dsc}}
\providecommand{\dtPG}{}\renewcommand{\dtPG}{dt}
\providecommand{\dtS}{}\renewcommand{\dtS}{dt_{\mathrm{S}}}
\providecommand{\ellv}{}\renewcommand{\ellv}{\ell}
\providecommand{\er}{}\renewcommand{\er}{e_r}
\providecommand{\et}{}\renewcommand{\et}{e_t}
\providecommand{\ev}{}\renewcommand{\ev}{\mathrm{ev}}
\providecommand{\gl}{}\renewcommand{\gl}{\mathfrak{gl}}
\providecommand{\hatV}{}\renewcommand{\hatV}{\widehat{V}}
\providecommand{\hb}{}\renewcommand{\hb}{\bar{h}}
\providecommand{\hf}{}\renewcommand{\hf}{\hat{f}}
\providecommand{\id}{}\renewcommand{\id}{\mathrm{id}}
\providecommand{\kpc}{}\renewcommand{\kpc}{\mathrm{kpc}}
\providecommand{\lP}{}\renewcommand{\lP}{\ell_{P}}
\providecommand{\mD}{}\renewcommand{\mD}{\mathbf{m}_{D}}
\providecommand{\mb}{}\renewcommand{\mb}{\bar m}
\providecommand{\mnu}{}\renewcommand{\mnu}{\Sigma m_{\nu}}
\providecommand{\nInf}{}\renewcommand{\nInf}{n_\infty}
\providecommand{\nO}{}\renewcommand{\nO}{n_0}
\providecommand{\ninf}{}\renewcommand{\ninf}{n_{\infty}}
\providecommand{\nz}{}\renewcommand{\nz}{n_{0}}
\providecommand{\proj}{}\renewcommand{\proj}{\operatorname{proj}}
\providecommand{\pv}{}\renewcommand{\pv}{\pi}
\providecommand{\rad}{}\renewcommand{\rad}{\mathrm{rad}}
\providecommand{\sgn}{}\renewcommand{\sgn}{\operatorname{sgn}}
\providecommand{\shA}{}\renewcommand{\shA}{\mathcal{A}}
\providecommand{\shC}{}\renewcommand{\shC}{\mathcal{C}}
\providecommand{\shE}{}\renewcommand{\shE}{\mathcal{E}}
\providecommand{\shJ}{}\renewcommand{\shJ}{\mathcal{J}}
\providecommand{\shK}{}\renewcommand{\shK}{\mathcal{K}}
\providecommand{\shL}{}\renewcommand{\shL}{L}
\providecommand{\shW}{}\renewcommand{\shW}{\mathcal{W}}
\providecommand{\sll}{}\renewcommand{\sll}{\mathfrak{sl}}
\providecommand{\so}{}\renewcommand{\so}{\mathfrak{so}}
\providecommand{\su}{}\renewcommand{\su}{\mathfrak{su}}
\providecommand{\supp}{}\renewcommand{\supp}{\operatorname{supp}}
\providecommand{\tr}{}\renewcommand{\tr}{\operatorname{tr}}
\providecommand{\uu}{}\renewcommand{\uu}{\mathfrak{u}}
\providecommand{\br}{}\renewcommand{\br}[1]{\langle #1\rangle}
\providecommand{\bra}{}\renewcommand{\bra}[1]{\langle #1|}
\providecommand{\file}{}\renewcommand{\file}[1]{\texttt{#1}}
\providecommand{\hp}{}\renewcommand{\hp}[2]{( #1 \,|\, #2 )}
\providecommand{\ip}{}\renewcommand{\ip}[2]{\langle #1, #2\rangle}
\providecommand{\ket}{}\renewcommand{\ket}[1]{|#1\rangle}
\providecommand{\rel}{}\renewcommand{\rel}[2]{D\bigl(#1 \,\|\, #2\bigr)}
\providecommand{\rev}{}\renewcommand{\rev}[1]{\widetilde{#1}}
\providecommand{\slashed}{}\renewcommand{\slashed}[1]{{#1}\mkern-10mu/}
\title{\bfseries The Matter in the Dark Towers}
\author{Leon Butler\\ \normalsize\href{mailto:lbutler.math@gmail.com}{\texttt{lbutler.math@gmail.com}}\\ \normalsize ORCID: \href{https://orcid.org/0009-0006-4449-5617}{0009-0006-4449-5617}}
\date{Generated by the Physics Graph, 20 November 2026 \textperiodcentered{} Transport Papers, T14c\\[3pt]\small Official version: \href{https://doi.org/10.5281/zenodo.24403592}{doi:10.5281/zenodo.24403592}}
\begin{document}
\maketitle
\begin{abstract}
\normalsize
The dark-matter paper T14 fixes how much matter each of five dark towers holds, not what that matter is. This note builds it with the Standard-Model paper's construction. A constituent is an element of the spin-irreducible even ideal of the tower's conformal algebra, and its scalars are the commutant of the spin action there. We prove that this field is the division algebra of the next residue: the reals in the three-dimensional towers, the complex numbers in this tower and the mirror, the quaternions in the five-dimensional towers. One identity of the reversion pairing then says where mass begins (two constituents in three and four dimensions, one in five); the momentum-preserving mixings are the field's unitary groups; and, under the series' identification premise and one transposition rule, the gauge algebras follow. The mirror carries the Standard Model again. Both three-dimensional towers carry the abelian $\so(2) \oplus \so(2)$: two copies of single-flavour electrodynamics in $2+1$ dimensions. Both five-dimensional towers carry $\uu(3)$, dark colour and one abelian charge, with vector-like quarks, leptons and a neutral singlet; with an external imaginary unit both would carry the Standard-Model algebra instead. The towers of each odd dimension differ only in discrete structure. Three consequences follow. Neither the unit-ball rule's premise E nor T14's particle model has a mechanism in any odd tower. A cold mirror is an unstated exception to T14's indifference premise. No configuration of the towers as built passes the cluster, radiation and mirror bounds together. The algebraic results are verified in exact arithmetic.
\end{abstract}
\newpage
\noindent \textbf{Authorship and disclosure. }The founding idea of the series is the author's: that division by zero in a Clifford algebra should be read not as a prohibition but as a transport into a higher-dimensional algebra, and that the transported division need not be accompanied by an inverse operation. The direction of investigation and the critical review of the results are likewise the author's. The mathematical development and the drafting of the text were carried out in extended collaboration with Claude Fable 5 (Anthropic), an AI system.

\noindent\fbox{\parbox{\dimexpr\linewidth-2\fboxsep-2\fboxrule}{\textbf{This is not the version of record.} The official version of this paper is the one deposited on Zenodo, \href{https://doi.org/10.5281/zenodo.24403592}{doi:10.5281/zenodo.24403592}. This PDF was generated by the Physics Graph on 20 November 2026 and assembled from approved parts of its store: the results, proofs and text of the paper as the store holds them, each result with its premises and machine checks. It is regenerated when a result it contains is corrected, so it may differ from the deposited paper. Since the deposited version 1.0: 28 results changed, 2 reworded (wording, notation or citations only) and 38 cited results added. Each revised or added result is marked in the margin (\textcolor{BrickRed}{corrected}, \textcolor{BurntOrange}{changed}, \textcolor{gray}{reworded}), and the last section, \emph{Changes since the deposited version}, lists them with the reason for each.}}
\newpage

\section{Introduction}

\phantomsection\label{c:note.dark-tower-matter.setting-setting-imports-premises}
Throughout, a tower is $V = \R^{p,q}$ with an orthonormal basis $\{e_{a}\}$, $Q(e_{a}) = e_{a}^{2} = \pm 1$, and $\Cl(V) = \Cl(p,q)$ its Clifford algebra, with even subalgebra $\Cl^{0}(V)$, reversion $x \mapsto \rev{x}$, grade involution $x \mapsto \hat{x}$ and grade projection $\ip{\cdot}_{k}$. In a causal signature ($\min(p,q) = 1$) the \emph{time direction} $e_{t}$ is the one basis vector of the minority sign, as in T14's convention \cite[Prop.~3.1]{ButlerT14}, and $Q(e_{t}) = \pm 1$ accordingly. The tower's conformal algebra is $\Cl(\hatV)$ with $\hatV = V \oplus \R^{1,1}$, that is $\Cl(p+1, q+1)$. The pseudoscalar of an algebra of dimension $n$ is $\omega = e_{1}e_{2}\cdots e_{n}$, central when $n$ is odd, with $Q(\omega) = (-1)^{n(n-1)/2}(-1)^{q}$.

The objects that the twistor paper T9 \cite{ButlerT9} and T12 describe in twistor language are, here, the even ideals of the conformal algebra and their elements; this note uses only the latter.

\begin{assembled}[Imported: the towers and their residues]\label{c:dark-tower-matter.imported-towers-their-types}
\changedmark{BurntOrange}{changed since the deposit}
The residue $r = p - q \bmod 8$ is constant along each tower of $(1,1)$ ascents. Under T14's Premise P3 the representatives are $(1,3)$ at $r = 6$, $(3,1)$ at $r = 2$, $(2,1)$ and $(1,2)$ at $r = 1, 7$, $(4,1)$ and $(1,4)$ at $r = 3, 5$, $(2,2)$ at $r = 0$ and $(4,0)$ or $(0,4)$ at $r = 4$; $r = 0$ and $r = 4$ are sterile. The odd residues carry no orientation bit and therefore no asymmetry. The towers interact only through the one metric (exact superselection). Sources: \cite[Prop.~3.1, Thm.~3.4, Thm.~3.10, Prop.~4.1]{ButlerT14}, \cite[Thm.~3.6]{ButlerT3}.
\end{assembled}
\begin{proof}
The residue $r = p - q \bmod 8$ is constant along each tower of $(1,1)$ ascents. Under T14's Premise P3 (\ref{c:dark-matter.p3-nearest-dimension-representatives}) the representatives are $(1,3)$ at $r = 6$, $(3,1)$ at $r = 2$, $(2,1)$ and $(1,2)$ at $r = 1, 7$, $(4,1)$ and $(1,4)$ at $r = 3, 5$, $(2,2)$ at $r = 0$ and $(4,0)$ or $(0,4)$ at $r = 4$; $r = 0$ and $r = 4$ are sterile. The odd residues carry no orientation bit and therefore no asymmetry. The towers interact only through the one metric (exact superselection). Sources: \cite[Prop.~3.1, Thm.~3.4, Thm.~3.10, Prop.~4.1]{ButlerT14} (\ref{c:dark-matter.superselected-hence-dark-exactly-stable}, \ref{c:dark-matter.flag-bit-exists-exactly-tower-dimension}, \ref{c:dark-matter.towers-are-residues-single-clock-filter}, \ref{c:dark-matter.type-residue-what-each-tower-can}), \cite[Thm.~3.6]{ButlerT3} (\ref{c:tower.rigidity-clock-frozen}).

\end{proof}

\begin{fact}[Imported: the structure of real Clifford algebras]\label{c:dark-tower-matter.imported-clifford-structure}
(i) For $q \geq 1$, $\Cl^{0}(p,q) \cong \Cl(p,q-1)$, by $e_{a} \mapsto e_{a}e_{n}$ with $Q(e_{n}) = -1$; for $p \geq 1$, $\Cl^{0}(p,q) \cong \Cl(q, p-1)$. (ii) $\Cl(p,q)$ is a matrix algebra over $D \in \{\R, \C, \HH\}$, or a sum of two, with $D$ fixed by $p - q \bmod 8$: $\R$ at $0, 1, 2$; $\C$ at $3, 7$; $\HH$ at $4, 5, 6$. (iii) A minimal left ideal $Af$ of a simple or semisimple algebra $A$, $f$ a primitive idempotent, has commutant $\mathrm{End}_{A}(Af) \cong (fAf)^{\mathrm{op}}$, a division algebra (Schur), acting by right multiplication. Sources: \cite{Chevalley1954, Lounesto2001, LawsonMichelsohn1989}.
\emph{Source:} C. Chevalley, The Algebraic Theory of Spinors, Columbia University Press (1954).; P. Lounesto, Clifford Algebras and Spinors, 2nd ed., Cambridge University Press (2001).; H. B. Lawson and M.-L. Michelsohn, Spin Geometry, Princeton University Press (1989)..
\end{fact}

\begin{assembled}[Imported: our tower's Standard-Model chain]\label{c:dark-tower-matter.imported-our-towers-chain}
In this tower: a single constituent has null momentum and two have the squared mass of T11's two-constituent formula \cite[Prop.~4.1]{ButlerT11}, \cite[Fact~2.1]{ButlerT12}; each constituent carries three ladder pairs \cite[Fact~2.2, Lem.~2.12]{ButlerT12}, the tower's rungs satisfying the canonical anticommutation relations \cite[Prop.~4.1]{ButlerT3}; the mixings of $n$ constituents preserving their total momentum and helicity form are $U(n)$ \cite[Prop.~3.1]{ButlerT12}; and the content, hypercharges, anomalies and Higgs representation are \cite[Thm.~4.1, Thm.~4.2, Prop.~5.3, Thm.~5.5, Prop.~9.22]{ButlerT12}.
\end{assembled}
\begin{proof}
In this tower: a single constituent has null momentum and two have the squared mass of T11's two-constituent formula \cite[Prop.~4.1]{ButlerT11} (\ref{c:mass.mass-symplectic-area-two-twistor-formula}), \cite[Fact~2.1]{ButlerT12}; each constituent carries three ladder pairs \cite[Fact~2.2, Lem.~2.12]{ButlerT12} (\ref{c:standard-model.pair-count-dimension}, \ref{c:standard-model.imported-pair-systems-two-chirality-halves}), the tower's rungs satisfying the canonical anticommutation relations \cite[Prop.~4.1]{ButlerT3} (\ref{c:tower.ladder-pairs-satisfy-car-vectors}); the mixings of $n$ constituents preserving their total momentum and helicity form are $U(n)$ \cite[Prop.~3.1]{ButlerT12} (\ref{c:standard-model.internal-symmetry-ladder}); and the content, hypercharges, anomalies and Higgs representation are \cite[Thm.~4.1, Thm.~4.2, Prop.~5.3, Thm.~5.5, Prop.~9.22]{ButlerT12} (\ref{c:standard-model.quarks-leptons-two-level-fermion-content}, \ref{c:standard-model.hypercharge-global-form-pair-counting-consistency}, \ref{c:standard-model.global-form-quotient-determined}, \ref{c:standard-model.ten-transported-ambient-space-doublets-are}, \ref{c:standard-model.gauge-algebra-minimal-massive-composite}).

\end{proof}

\begin{assembled}[Imported: the unit-ball shares]\label{c:dark-tower-matter.imported-unit-ball-rule}
\changedmark{BurntOrange}{changed since the deposit}
T14's unit-ball rule rests on two premises: E, one length $d$ common to every populated tower at which each holds the same mass per ball of radius $d$ of its own space, and R3, a three-dimensional tower's sheets stacked one $d$ apart and a five-dimensional tower's matter counted through one $d$ of its extra direction. Under E and R3 the shares are $1$, $2 \times 4/3$, $2 \times 8/(3\pi)$, with total $(11\pi + 16)/(3\pi)$ \cite[Prem.~7.2, Prem.~7.4, Prop.~7.6(i)]{ButlerT14}. T14's particle model, R1 (every quantum of every populated tower has the same mass) and R2 (one quantum per unit ball of each tower's own space, so a common Wigner-Seitz radius $d$), implies E and is not used by the rule \cite[Prem.~7.3, Prop.~7.6(ii)]{ButlerT14}; in it $d$ is this world's baryon spacing, $d_{b} \approx 0.98$ m, and the mass the nucleon's \cite[Rem.~7.9, Cor.~7.26]{ButlerT14}. E alone fixes neither. As fractions of the dark matter: mirror $18.6$ percent, each three-dimensional tower $24.9$, each five-dimensional tower $15.8$.
\end{assembled}
\begin{proof}
Imported from T14. Premise E (\ref{c:dark-matter.e-equal-mass-per-unit-ball}) (one length $d$ common to every populated tower, at which each holds the same mass per ball of radius $d$ of its own space) and R3 (\ref{c:dark-matter.r3-one-cell-per-unshared-direction}) (a three-dimensional tower's sheets stacked one $d$ apart; a five-dimensional tower's matter counted through one $d$ of its extra direction) give shares $1$, $2 \times 4/3$, $2 \times 8/(3\pi)$ and the total $(11\pi + 16)/(3\pi)$ \cite[Prem.~7.2, Prem.~7.4, Prop.~7.6(i)]{ButlerT14} (\ref{c:dark-matter.unit-ball-rule}). The particle model, R1 (\ref{c:dark-matter.r1-equal-mass-per-quantum}) and R2 (\ref{c:dark-matter.r2-one-quantum-per-unit-ball}), implies E and is not used by the rule; it supplies the values $d = d_{b}$ and the nucleon's mass \cite[Prem.~7.3, Prop.~7.6(ii)]{ButlerT14}. As fractions of the dark matter: mirror $18.6$ percent, each three-dimensional tower $24.9$, each five-dimensional tower $15.8$.

\end{proof}

\begin{fact}[Imported: the cluster merger simulations]\label{c:dark-tower-matter.imported-cluster-simulations}
\changedmark{BurntOrange}{changed since the deposit}
In the idealised Bullet-like merger runs of T14b \cite[Obs.~7.1, Table~3]{ButlerT14b}: (i) with only the mirror colliding ($18.6$ percent of the dark matter, as gas) the subcluster keeps a separate lensing peak on its galaxies and a mass-to-light change of $0.82$ against the observed $0.84 \pm 0.07$; (ii) with the mirror gas and one three-dimensional tower's share of colliding fluid ($43.5$ percent colliding) the subcluster dissolves, and with more it dissolves too; (iii) with the mirror collisionless and $24.9$ or $31.6$ percent of colliding fluid, the subcluster keeps a weaker peak but only $0.55$ or $0.44$ of its central mass: the colliding fluid alone costs about half of it. In (i) the subcluster's own mass is $0.57$ of the cold-dark-matter run's, but mass and light leave together and the mass-to-light change, the quantity the observational bound is inferred from, is consistent. No run lies between $18.6$ and $24.9$ percent colliding. The colliding components there are fluids.
\emph{Source:} L. Butler, Simulating the Dark Towers: Methods, Assumptions and Tests of the Unit-Ball Composition, Transport Papers T14b (2026), Obs. 7.1, Table 3.
\end{fact}

\phantomsection\label{c:note.dark-tower-matter.setting-setting-imports-premises-2}
Four statements are carried as premises. H1 and H2 are this paper's names for two of T12's ingredients used off dimension four; Hypothesis (S) is T12's; U is this paper's transposition rule.

\section*{What we assume}
The results below are taken as given in this paper; each result that rests on one says so.
\begin{postulate}[E: equal mass per unit ball]\label{c:dark-matter.e-equal-mass-per-unit-ball}
\emph{Premise (E), equal mass per unit ball.} There is one length $d$, common to all populated towers, such that every populated tower holds the same mass $\mu$ per ball of radius $d$ of its own space. With $\rho_{r}$ the mass per unit volume of the $D_{r}$-dimensional space of the tower of residue $r$, at its representative under Premise \ref{c:dark-matter.p3-nearest-dimension-representatives}, and $V_{D}$ the volume of the unit ball $\B^{D}$,
\[
\rho_{r}\,V_{D_{r}}\,d^{D_{r}} \;=\; \mu \qquad \text{for every populated } r .
\]
The unit ball of the unit-ball rule is the ball of radius $d$.

\emph{What it replaces.} Premises R1 (\ref{c:dark-matter.r1-equal-mass-per-quantum}, equal mass per quantum) and R2 (\ref{c:dark-matter.r2-one-quantum-per-unit-ball}, one quantum per unit ball), taken in the length $d$, give E with $\mu$ the mass of a quantum, and E gives neither: it says nothing about what a tower's quanta weigh or how many there are (Theorem \emph{Dark-to-baryon total uses R1, R2 only as mass per unit ball at a fixed common thickness}(i)). R1 and R2 are one model of E, the particle model. By the decision of 19 November 2026, E and not R1 and R2 is the premise of the dark-to-baryon total (Proposition \ref{c:dark-matter.unit-ball-rule}), with the thickness of Premise R3 (\ref{c:dark-matter.r3-one-cell-per-unshared-direction}) read in the same length $d$. R1 and R2 remain premises of the stored results that use a quantum's mass or number by itself.

\emph{What it does not fix.} E does not fix $d$ independently of the towers' densities, which are not observed, and the total does not depend on it: once E and R3 hold, $d$ is $\rho_{1}V_{2}/(\rho_{6}V_{3})$, the density of a three-dimensional tower against ours (Proposition \ref{c:dark-matter.unit-ball-rule}(iii)). In particular $d$ is not tied to the spacing of this world's baryons. That tie belongs to the particle model, R1 and R2 together: E at $d$ with R2 in another common length is consistent, the quanta then having unequal masses. The stored results that give the counted thickness a value, today's baryon spacing of about a metre or its value at the deposit, rest on R1 and R2 for that. The democracy premise (\emph{D: democracy of the total mass} [T14, Premise 7.19]) is E with $d$ the de Sitter radius, given the stored reading of the sphere form, and the particle model is E with $d$ the common cell radius (Theorem \emph{Dark-to-baryon total uses R1, R2 only as mass per unit ball at a fixed common thickness}(i), (ii)). E can hold in one length only (Proposition \ref{c:dark-matter.unit-ball-rule}(iv)). It presupposes that the towers it names are populated, which for the four odd towers is Premise \emph{Pop: the four odd towers are populated}. Whether one length serves at more than one epoch is settled as for R2 (Theorem \emph{The unit-ball clauses hold at one epoch unless the hidden direction co-expands}(v)): only if the five-dimensional towers' hidden direction co-expands with ours, and otherwise E is a statement about the deposit.

\emph{How much it says.} For the mirror, whose space has the dimension of ours, E is equal density, which the population count gives separately on the proton-only branch (Proposition \emph{The population count: one share to half a percent, and a symmetric remainder} [T14, Prop 4.5]; on the burning branch the count gives $0.9948$ and E fails for the mirror by about one part in two hundred, Theorem \emph{Dark-to-baryon total uses R1, R2 only as mass per unit ball at a fixed common thickness}(iii)). E with R3 holds for some $d$ exactly when the closed form's shares hold tower by tower and the two thicknesses are equal (Proposition \ref{c:dark-matter.unit-ball-rule}(iii)). So E is the fitted relation stated as one law: it adds no support to the fit that is independent of the fitted number, and the series supplies no mechanism for it.

\emph{Refutation condition.} Observation can refute E with R3 only through the closed form: by a dark-to-baryon ratio that excludes $(11\pi + 16)/3\pi$, or by a composition of the dark matter other than the shares $V_{3}/V_{D}$. E with R3 says more than the total does, that the shares hold tower by tower and that the two thicknesses are equal (Proposition \ref{c:dark-matter.unit-ball-rule}(iii)), and that further content is not observed.
\end{postulate}

\begin{postulate}[P3: nearest-dimension representatives]\label{c:dark-matter.p3-nearest-dimension-representatives}
Each tower is represented at the dimension nearest four; where two are equally near (the odd residues, three and five), the lower one with a single time, else the higher. \emph{Support:} P1's one metric read as minimality. Structural premise; Proposition \ref{c:dark-matter.towers-are-residues-single-clock-filter}(iii) records the other windows. \emph{Note.} Premise (W) of \emph{The frozen point's complement as window: lowest causal representatives}, that every tower's spacetime is, up to sign, a nondegenerate subspace of the frozen point's complement, is recorded as a costed alternative that restates this premise's window (with the single-clock filter it would replace this premise and the residue-$4$ half of the costed single-ideal clause of Observation \emph{The status of the rule} [T14, Obs 7.31](vi) by one hypothesis; at residue $0$ it admits the causal $(1,1)$ where this premise has $(2,2)$); it was found with the fit known and is not adopted. This premise is unchanged. \emph{Note (by residue).} At the odd residues this premise takes the lowest causal member, $(2,1)$, $(1,2)$ at residues $1$, $7$ and $(4,1)$, $(1,4)$ at residues $3$, $5$. At residues $1$ and $7$ that choice follows from Premise U (\ref{c:dark-tower-matter.u-transposition-rule}): the constituent field is $\R$ in every causal member and the stored pair count exceeds the real Witt index $2$ above dimension three, so U admits only $(2,1)$ and $(1,2)$. At residues $3$ and $5$ it is independent of the descent: a model with those towers at $(12,1)$ and $(1,12)$ satisfies every premise of the window-sensitive results other than this one, Premise U included, so there the window is an input (Theorem \emph{The five-dimensional towers' dimension is free; Premise U fixes the three-dimensional ones}(ii), (iii)).
\end{postulate}

\begin{postulate}[R1: equal mass per quantum]\label{c:dark-matter.r1-equal-mass-per-quantum}
The quanta of every populated tower carry the same mass $m$.
\end{postulate}

\begin{postulate}[R2: one quantum per unit ball]\label{c:dark-matter.r2-one-quantum-per-unit-ball}
A populated tower whose space has $D$ dimensions (under P3, \ref{c:dark-matter.p3-nearest-dimension-representatives}, $2$, $3$ or $4$) holds one quantum per unit ball $\B^{D}$ of that space, in one unit of length common to all towers: every tower has the same Wigner-Seitz radius. The rule presupposes that all six towers are populated, which for the odd four is the inference of Observation \emph{Which towers are populated} [T14, Obs 4.10].
\end{postulate}

\begin{postulate}[R3: one cell per unshared direction]\label{c:dark-matter.r3-one-cell-per-unshared-direction}
In the sourcing rule, a tower's matter counts through exactly one unit of length of each direction its neighbour lacks, on one side: a five-dimensional tower's matter through one unit of its extra direction, and a three-dimensional tower's sheets stacked one unit apart along the direction it lacks. the section ``The thickness: a contact mechanism'' gives a mechanism that returns this value and costs it.

\emph{The unit.} The unit is the length $d$ of Premise E (\ref{c:dark-matter.e-equal-mass-per-unit-ball}): one length common to all towers, the radius of the ball in which E counts mass, and independent of how a tower's mass is divided into quanta. This is the fixed-common-length reading, adopted by the decision of 19 November 2026; it replaces the reading of 10 November 2026, in which the unit was each tower's own cell radius, the Wigner-Seitz radius of its own filling. Under R1 and R2 (\ref{c:dark-matter.r1-equal-mass-per-quantum}, \ref{c:dark-matter.r2-one-quantum-per-unit-ball}), E holds in R2's common length, which is then $d$ (Proposition \ref{c:dark-matter.unit-ball-rule}(iv)), every tower's cell radius is $d$ and the two readings coincide; under R2 alone they coincide only if R2 is taken in the length $d$. They differ when R2 fails: with a tower's own cell radius $a_{r}$ as the unit the shares are $(m_{r}/m_{6})(V_{3}/V_{D_{r}})(a_{6}/a_{r})^{3}$, which Premise E does not fix, and with the common length they are $V_{3}/V_{D_{r}}$ under E alone (Theorem \emph{Dark-to-baryon total uses R1, R2 only as mass per unit ball at a fixed common thickness}(i)). The title keeps the word cell: the cell is the unit ball of Premise E. The premise takes its unit from E, so where E is not assumed the unit has to be supplied in another way: by R1 with R2, whose common cell radius is then E's length, or by the democracy premise (\emph{D: democracy of the total mass} [T14, Premise 7.19]), whose length is the de Sitter radius.
\end{postulate}

\begin{postulate}[Premise H1: a constituent's momentum]\label{c:dark-tower-matter.h1-momentum}
\changedmark{BurntOrange}{changed since the deposit}
A constituent has a spinor part $\psi$ in the tower's spinor ideal $S$ (Definition \ref{c:obj.dark-tower-matter.ideals-fields-momentum}), carried by a map $T \to S$ that intertwines the action of the tower's spin group, so that the field of $T$ acts on $S$ through it; its momentum is the vector part of $\psi e_{t}\rev{\psi}$. In dimension four this is T12's Fact 2.1 in the even algebra; in dimensions three and five it is a premise.

\emph{Note (what of H1 is forced).} In the causal towers of dimension three to five the quadratic vectors of one constituent's spinor part that are covariant under the connected spin group form one real line, spanned by this momentum (Theorem \emph{Causal towers of dimension 3 to 5: j the only covariant quadratic vector of a constituent}). So, given a spinor part and a momentum that is quadratic and covariant, the formula is forced up to a real factor. H1 is weakened by that result, not retired: what stays premised in dimensions three and five is that a constituent has a spinor part in $S$ carried by an intertwining map, that its momentum is a quadratic covariant vector of it, the normalisation, and the additivity of momenta over constituents, which covariance does not force.
\end{postulate}

\begin{postulate}[Premise H2: gauge algebras as hierarchy automorphisms]\label{c:dark-tower-matter.h2-identification}
The gauge algebra of a tower's matter is the automorphism algebra of its minimal massive hierarchy, constituent slots tensor pair slots, as T12's Theorem 4.1 identifies it in dimension four, at observation level \cite[Thm.~4.1, Prop.~4.3]{ButlerT12}.
\end{postulate}

\begin{postulate}[Hypothesis (S) carried to every tower]\label{c:dark-tower-matter.hypothesis-s-slot-statistics}
The constituent slots are fermionic modes on the same footing as the pairs, so that the composite's state space is the exterior algebra of slot and pair modes \cite[Sec.~5]{ButlerT12}. Used only where stated.
\end{postulate}

\begin{postulate}[Premise U: the transposition rule (a definition)]\label{c:dark-tower-matter.u-transposition-rule}
\changedmark{BurntOrange}{changed since the deposit}
(U1) A constituent is an element of the spin-irreducible even ideal of the conformal algebra, as T12 defines this tower's, and its field is the commutant of the spin action on it. (U2) The hierarchy's first factor counts constituents at mass onset, and its slot group is the momentum-preserving group of that many constituents; the pair group is the group of mixings of the ladder pairs that preserves their canonical relations and the field's positive norm, the unitary group over the pairs' scalars ($\R$, or $\C$ where a complex unit builds the pairs), as T12 takes $U(m) \times U(k)$ \cite[Thm.~4.1]{ButlerT12}. The norm clause is part of the definition: no positive form on the pair space is derived, and without it the group preserving the canonical relations is $GL(2,\R)$; the compactness of every gauge group stated under U rests on it. (U3) Every complex unit used to build a ladder pair is an element of the constituent's field that the real algebra supplies; no external imaginary unit is adjoined.

\emph{Type: a definition.} The three clauses define what the series calls a constituent in another tower, its field, the slot and pair groups, and which complex units are admitted; they assert nothing about nature that could be derived or refuted, and can only be replaced by another definition (Remark \ref{c:note.dark-tower-matter.branches-break-u-what-they-give} lists what other choices give). The definition has consequences through what it excludes: with it, the lowest-causal choice of Premise \ref{c:dark-matter.p3-nearest-dimension-representatives} at residues $1$ and $7$ follows. Results stated under Premise U hold for the objects so defined; the entity keeps the status of a premise so that those results show the choice among their conditions.
\end{postulate}

\section{Setting, imports and premises}

\phantomsection\label{c:note.dark-tower-matter.setting-constituents-momentum-where-mass-begins}
The commutant of the spin action is the commutant of the even algebra, since the bivectors generate $\Cl^{0}$ as an algebra; so U1's field and the field of this definition agree. The constructions are explicit: $e_{t}e_{a}$ squares to $+1$, and $\Cl^{0}(V)f$ has real dimension $2$, $4$, $8$ in dimensions three, four, five, which is minimal (Theorem \ref{c:dark-tower-matter.constituents-field}); $\hf$ is a product of two or three commuting factors $\tfrac{1}{2}(1 + b)$ with $b$ an even blade of square $+1$. The script builds both by a search over commuting even blades (Part A).

\begin{theorem}[The constituent's field in each tower]\label{c:dark-tower-matter.constituents-field}
\changedmark{BurntOrange}{changed since the deposit}
In tower $r$ the fields $\F(T)$ and $\F(S)$ are equal, and equal to the division algebra of residue $r + 1$:
\begin{center}\small
\begin{tabular}{@{}lccccccccc@{}}
\toprule
$r$ & $6$ & $2$ & $1$ & $7$ & $3$ & $5$ & $0$ & $4$ \\
tower & $(1,3)$ & $(3,1)$ & $(2,1)$ & $(1,2)$ & $(4,1)$ & $(1,4)$ & $(2,2)$ & $(4,0)$, $(0,4)$ \\
conformal & $(2,4)$ & $(4,2)$ & $(3,2)$ & $(2,3)$ & $(5,2)$ & $(2,5)$ & $(3,3)$ & $(5,1)$, $(1,5)$ \\
field & $\C$ & $\C$ & $\R$ & $\R$ & $\HH$ & $\HH$ & $\R$ & $\HH$ \\
$\dim_{\R} T$, $\dim_{\R} S$ & $8, 4$ & $8, 4$ & $4, 2$ & $4, 2$ & $16, 8$ & $16, 8$ & $4, 2$ & $8, 4$ \\
\bottomrule
\end{tabular}
\end{center}
In particular the field is not a function of the conformal pseudoscalar's square: $Q(\omega) = -1$ at $r = 7$ (field $\R$) and at $r = 3$ (field $\HH$).
\end{theorem}
\begin{proof}
\textit{The two even algebras.}\enspace By Fact \ref{c:dark-tower-matter.imported-clifford-structure}(i),
\begin{align*}
\Cl^{0}(\hatV) &= \Cl^{0}(p+1, q+1) \cong \Cl(p+1, q), \\
\Cl^{0}(V) &= \Cl^{0}(p, q) \cong \Cl(p, q-1),
\end{align*}
and both have residue $(p+1) - q = p - (q-1) = r + 1$. The one tower with $q = 0$ is $(4,0)$: there $\Cl^{0}(4,0) \cong \Cl(0,3)$, of residue $-3 \equiv 5 = r + 1$, and $(0,4)$ gives the same algebra.



\emph{The commutant.}  By Fact \ref{c:dark-tower-matter.imported-clifford-structure}(ii) each algebra is a matrix algebra over the division algebra $D(r+1)$, or a sum of two copies of one, and by Fact \ref{c:dark-tower-matter.imported-clifford-structure}(iii) the commutant of a minimal left ideal is $D(r+1)$. Reading $D$ off for each tower:
\[
\begin{array}{c|cccccccc}
r & 6 & 2 & 1 & 7 & 3 & 5 & 0 & 4 \\
r + 1 & 7 & 3 & 2 & 0 & 4 & 6 & 1 & 5 \\ \hline
D(r+1) & \C & \C & \R & \R & \HH & \HH & \R & \HH
\end{array}
\]



\emph{Checks.}  The script computes every commutant from scratch on exact matrix representations of the ideals and classifies it (checks A1, A2), with the dimensions (A3) and the pseudoscalars (A4, A5).

\end{proof}

\section{Constituents, momentum and where mass begins}

\begin{definition}[Ideals, fields and momentum]\label{c:obj.dark-tower-matter.ideals-fields-momentum}
Let $\hf$ be a primitive idempotent of $\Cl^{0}(\hatV)$, and in a causal tower let $f = \tfrac{1}{2}(1 + e_{t}e_{a})$ for a space direction $e_{a}$, which is a primitive idempotent of $\Cl^{0}(V)$ in dimensions three to five (Theorem \ref{c:dark-tower-matter.constituents-field}); in the sterile towers $f$ is any primitive idempotent of $\Cl^{0}(V)$. The \emph{constituent ideal} is $T = \Cl^{0}(\hatV)\hf$ and the \emph{spinor ideal} is $S = \Cl^{0}(V)f$. Their \emph{fields} are the commutants $\F(T) = \mathrm{End}_{\Cl^{0}(\hatV)}(T)$ and $\F(S) = \mathrm{End}_{\Cl^{0}(V)}(S)$, the latter realised as right multiplication by the corner algebra $f\Cl^{0}(V)f$. In a causal tower the \emph{adjoint} is $x^{\dagger} = e_{t}\rev{x}e_{t}^{-1}$, and the \emph{momentum} of $\psi \in S$ is
\[
j(\psi) \;=\; \ip{\psi e_{t} \rev{\psi}}_{1}.
\]
The total momentum of constituents $\psi_{1}, \ldots, \psi_{n}$ is $P = \sum_{i} j(\psi_{i})$, and its \emph{squared mass} is $m^{2} = Q(e_{t})\,Q(P)$.
\end{definition}

\begin{remark}[T14's types are a different commutant]\label{c:note.dark-tower-matter.t14s-types-are-different-commutant}
T14's Theorem 3.4 assigns real type to $r = 1, 2$, complex type with an intrinsic $U(1)$ to $r = 3, 7$, and quaternionic type to $r = 4, 5, 6$ \cite[Thm.~3.4]{ButlerT14}. That is the commutant of the full Clifford algebra $\Cl(V)$ on its irreducible module; Theorem \ref{c:dark-tower-matter.constituents-field} is the commutant of the spin action on the constituent ideal. Both are correct and they answer different questions. T14's intrinsic $U(1)$ at $r = 3$ and $r = 7$ is the central pseudoscalar, which reappears here as structure the tower carries beside its constituent field: at $r = 7$ as a unit joining two real constituents (Proposition \ref{c:dark-tower-matter.residue-seven-unit}), at $r = 3$ as the unit $I_{7}$ inside the quaternions (Theorem \ref{c:dark-tower-matter.third-pairs-unit-five-dimensions}).
\end{remark}

\begin{lemma}[The adjoint and the norm of the field]\label{c:dark-tower-matter.adjoint-norm}
\changedmark{BurntOrange}{changed since the deposit}
In a causal tower:
\begin{enumerate}
\item[(i)] $e_{K}e_{K}^{\dagger} = Q(e_{t})^{|K|}$ for every basis blade $e_{K}$, so $e_{K}e_{K}^{\dagger} = 1$ on every even blade and, for even $x$, $\ip{xx^{\dagger}}_{0}$ is the sum of the squares of $x$'s coordinates;
\item[(ii)] $f^{\dagger} = f$, and for $x \in f\Cl^{0}(V)f$,
\[
x^{\dagger}x \;=\; xx^{\dagger} \;=\; N(x)\,f, \qquad N(x) = 2\ip{xx^{\dagger}}_{0} \geq 0,
\]
with $N(x) = 0$ only at $x = 0$;
\item[(iii)] for $\psi \in S$ and $x \in f\Cl^{0}(V)f$,
\[
j(\psi x) \;=\; N(x)\, j(\psi).
\]
\end{enumerate}
\end{lemma}
\begin{proof}
\textit{(i) The blades.}\enspace Every direction other than $e_{t}$ has square $-Q(e_{t})$. There are two cases.

If $t \notin K$, moving $e_{t}$ through the $|K|$ vectors of $\rev{e_{K}}$ costs one sign each, and a blade times its reverse is the product of its squares:
\[
e_{K}^{\dagger} = e_{t}\rev{e_{K}}e_{t}^{-1} = (-1)^{|K|}\rev{e_{K}}, \qquad e_{K}\rev{e_{K}} = \bigl(-Q(e_{t})\bigr)^{|K|},
\]
so $e_{K}e_{K}^{\dagger} = (-1)^{|K|}\bigl(-Q(e_{t})\bigr)^{|K|} = Q(e_{t})^{|K|}$.

If $e_{K} = e_{t}e_{K'}$, then $\rev{e_{K}} = \rev{e_{K'}}e_{t}$, and
\begin{align*}
e_{K}^{\dagger} &= e_{t}\rev{e_{K'}} = (-1)^{|K'|}\,\rev{e_{K'}}\,e_{t}, \\
e_{K}e_{K}^{\dagger} &= (-1)^{|K'|}\,e_{t}\bigl(e_{K'}\rev{e_{K'}}\bigr)e_{t} = (-1)^{|K'|}\bigl(-Q(e_{t})\bigr)^{|K'|}\,Q(e_{t}) = Q(e_{t})^{|K|}.
\end{align*}
For $K \neq L$, $e_{K}e_{L}^{\dagger}$ is a multiple of $e_{K \triangle L}$, which has no scalar part. Hence for even $x = \sum_{K}x_{K}e_{K}$,
\[
\ip{xx^{\dagger}}_{0} \;=\; \sum_{K,L}x_{K}x_{L}\,\ip{e_{K}e_{L}^{\dagger}}_{0} \;=\; \sum_{K}x_{K}^{2} \qquad \text{(check B0)}.
\]



\emph{(ii) The corner and its conjugation.}  With $f = \tfrac{1}{2}(1 + e_{t}e_{a})$ and $\rev{(e_{t}e_{a})} = e_{a}e_{t}$,
\[
f^{\dagger} = \tfrac{1}{2}\bigl(1 + e_{t}\,e_{a}e_{t}\,e_{t}^{-1}\bigr) = \tfrac{1}{2}(1 + e_{t}e_{a}) = f.
\]
The corner $\F = f\Cl^{0}f$ is a division algebra $\R$, $\C$ or $\HH$ with unit $f$ (Fact \ref{c:dark-tower-matter.imported-clifford-structure}(iii)). It is closed under $\dagger$, since $(fxf)^{\dagger} = fx^{\dagger}f$, and $\dagger$ is an anti-involution of it with $\ip{xx^{\dagger}}_{0} > 0$ for $x \neq 0$ by (i).

The standard conjugation is the only anti-involution with that positivity. Every other candidate has an element with $xx^{\dagger} = -f$, whose scalar part is $-\tfrac{1}{2}$:
\begin{itemize}
\item on $\C$ the other candidate is the identity, and a unit square root $x$ of $-f$ gives $xx^{\dagger} = xx = -f$;
\item on $\HH$ every other anti-involution is $x \mapsto u\bar{x}u^{-1}$ for a unit imaginary $u$, and a unit imaginary $v$ anticommuting with $u$ gives $v\,u\bar{v}u^{-1} = v^{2} = -f$.
\end{itemize}
So $\dagger$ is the standard conjugation, and
\[
x^{\dagger}x \;=\; xx^{\dagger} \;=\; N(x)\,f
\]
with $N$ the standard norm; the script verifies this as a polynomial identity in $x$'s coordinates (check B1). Taking scalar parts, with $\ip{f}_{0} = \tfrac{1}{2}$,
\[
\ip{xx^{\dagger}}_{0} \;=\; N(x)\,\ip{f}_{0} \;=\; \tfrac{1}{2}\,N(x),
\]
which is positive definite by (i), since $x$ is even (B3).



\emph{(iii) The norm law.}  By the definition of the adjoint, $e_{t}\rev{x} = x^{\dagger}e_{t}$, and $\psi f = \psi$ for $\psi \in S$. So
\begin{align*}
\psi x\,e_{t}\,\rev{(\psi x)} &= \psi\,\bigl(x e_{t}\rev{x}\bigr)\,\rev{\psi} \\
&= \psi\,\bigl(xx^{\dagger}\bigr)\,e_{t}\rev{\psi} \\
&= N(x)\,\psi f\,e_{t}\rev{\psi} \\
&= N(x)\,\psi e_{t}\rev{\psi},
\end{align*}
and the grade-one part is $j(\psi x) = N(x)\,j(\psi)$ (B2).

\end{proof}

\begin{theorem}[Where mass begins in each causal tower]\label{c:dark-tower-matter.where-mass-begins}
\changedmark{BurntOrange}{changed since the deposit}
Let $\psi, \psi_{1}, \psi_{2} \in S$ in a causal tower of dimension $n$.
\begin{enumerate}
\item[(i)] For $n = 3$ and $n = 4$ the reversion pairing is alternating, $\rev{\psi}\psi = 0$, every single constituent's momentum is null, and
\[
m^{2}\bigl(j(\psi_{1}) + j(\psi_{2})\bigr) \;=\; 2\,\ip{xx^{\dagger}}_{0}, \qquad x = \rev{\psi}_{1}\psi_{2}.
\]
Mass begins at two constituents; it vanishes exactly when $\rev{\psi}_{1}\psi_{2} = 0$.
\item[(ii)] For $n = 5$ the reversion pairing is not alternating, and
\[
m^{2}\bigl(j(\psi)\bigr) \;=\; \tfrac{1}{2}\,\ip{xx^{\dagger}}_{0}, \qquad x = \rev{\psi}\psi,
\]
which vanishes on a nonzero proper cone in $S$. Mass begins at one constituent.
\end{enumerate}
In every causal tower the total momentum of any number of constituents is timelike or null, never spacelike.
\end{theorem}
\begin{proof}
\textit{The pairing as a sandwich.}\enspace Write $\psi = \phi f$ with $\phi \in \Cl^{0}(V)$ and $f = \tfrac{1}{2}(1 + b)$, $b$ a bivector of square $+1$. Then $\rev{b} = -b$, so $\rev{f} = 1 - f$ and
\[
\rev{\psi}\psi \;=\; (1 - f)\;\rev{\phi}\phi\;f.
\]
The element $\rev{\phi}\phi$ is even and reversion-invariant, so it has grades $0$ and $4$ only:
\[
\rev{\phi}\phi = s + F, \qquad s = \ip{\rev{\phi}\phi}_{0}, \quad F = \ip{\rev{\phi}\phi}_{4}.
\]
Since $(1 - f)f = 0$, every part of $\rev{\phi}\phi$ that commutes with $b$ drops out of the sandwich:
\[
(1 - f)\,c\,f \;=\; c\,(1 - f)f \;=\; 0 \qquad \text{whenever } cb = bc.
\]



\emph{Dimensions three and four.}  In dimension three $F = 0$. In dimension four $F = uI$ is a multiple of the pseudoscalar, which commutes with every even element and in particular with $b$. In both cases all of $\rev{\phi}\phi$ commutes with $b$, so $\rev{\psi}\psi = 0$. This is the alternation in (i).



\emph{Dimension five.}  Here $F$ is a general four-vector, which need not commute with $b$, and $(1 - f)\,F f \neq 0$ for suitable $\phi$, so the pairing is not alternating.



\emph{The mass laws.}  The two displayed laws are identities between polynomials of degree four in the coordinates of $\psi_{1}, \psi_{2}$ (respectively $\psi$) on a basis of $S$. The script verifies them as polynomial identities over $\mathbb{Q}$, not on samples, in all six causal towers (checks B4, B5), together with the nullity of $j(\psi)$ in dimensions three and four (B4); it exhibits a nonzero $\psi$ with $\rev{\psi}\psi = 0$ in each five-dimensional tower (B6). The right-hand sides are nonnegative by Lemma \ref{c:dark-tower-matter.adjoint-norm}(i), so each single momentum (and in dimensions three and four each pair sum) is timelike or null.



\emph{Any number of constituents.}  For every $\psi \in S$,
\begin{align*}
Q(e_{t})\,\bigl(j(\psi)\cdot e_{t}\bigr) &= Q(e_{t})\,\ip{\psi e_{t}\rev{\psi}\,e_{t}}_{0} \\
&= \ip{\psi\,e_{t}\rev{\psi}e_{t}^{-1}}_{0} \\
&= \ip{\psi\psi^{\dagger}}_{0} \;\geq\; 0 \qquad \text{(B7)}.
\end{align*}
So every constituent's momentum is causal and lies in the same closed half-space bounded by $e_{t}^{\perp}$, hence in the same half of the causal cone, and that half is closed under addition.

\end{proof}

\begin{remark}[The four-dimensional case is T11's formula]\label{c:note.dark-tower-matter.four-dimensional-case-t11s-formula}
In this tower $x = \rev{\psi}_{1}\psi_{2}$ lies in the two-dimensional space $(1-f)\Cl^{0}f$, a copy of $\C$, and Theorem \ref{c:dark-tower-matter.where-mass-begins}(i) is T11's two-constituent mass formula \cite[Prop.~4.1]{ButlerT11} written in the even algebra, the pairing $x$ playing the role of the bracket. In five dimensions the law is the classical Fierz identity of a single Dirac spinor, which in geometric algebra says that the four-vector part of $\rev{\phi}\phi$ survives the idempotent.
\end{remark}

\begin{theorem}[The momentum-preserving mixings of constituents]\label{c:dark-tower-matter.ladder}
\changedmark{BurntOrange}{changed since the deposit}
Let $\F = f\Cl^{0}(V)f$ act on $S$ on the right and on $S^{n}$ by $\psi'_{i} = \sum_{k}\psi_{k}U_{ki}$ for $U \in M_{n}(\F)$. Then $U$ preserves the total momentum of every $n$-tuple if and only if $UU^{\dagger} = 1$, where $(U^{\dagger})_{ik} = (U_{ki})^{\dagger}$. The momentum-preserving group is $O(n)$, $U(n)$ or $Sp(n)$ as $\F = \R$, $\C$, $\HH$: $O(n)$ in the three-dimensional towers, $U(n)$ in this tower and the mirror, $Sp(n)$ in the five-dimensional towers.
\end{theorem}
\begin{proof}
\textit{Unitary mixings preserve the momentum.}\enspace Using $e_{t}\rev{y} = y^{\dagger}e_{t}$ on each entry of $U$,
\begin{align*}
\sum_{i}\psi'_{i}e_{t}\rev{\psi'_{i}} &= \sum_{i}\sum_{k,l}\psi_{k}U_{ki}\,e_{t}\,\rev{U_{li}}\,\rev{\psi_{l}} \\
&= \sum_{k,l}\psi_{k}\Bigl(\sum_{i}U_{ki}U_{li}^{\dagger}\Bigr)e_{t}\rev{\psi_{l}}.
\end{align*}
If $UU^{\dagger} = 1$, that is $\sum_{i}U_{ki}U_{li}^{\dagger} = \delta_{kl}f$, the right-hand side is $\sum_{k}\psi_{k}e_{t}\rev{\psi_{k}}$.



\emph{Only unitary mixings do.}  Take $\psi_{k} = \psi a_{k}$ for one $\psi \in S$ with $j(\psi) \neq 0$ and a row $a = (a_{k}) \in \F^{n}$. Then $\psi'_{i} = \psi\,(aU)_{i}$, and Lemma \ref{c:dark-tower-matter.adjoint-norm}(iii) gives the total momenta before and after as
\[
P = \Bigl(\sum_{k}N(a_{k})\Bigr)\,j(\psi), \qquad P' = \Bigl(\sum_{i}N\bigl((aU)_{i}\bigr)\Bigr)\,j(\psi).
\]
Their equality for every $a$ says that $U$ preserves the positive definite form $\sum_{k}N(a_{k})$ on $\F^{n}$, whose polarisation is the standard $\F$-Hermitian form; so $UU^{\dagger} = 1$.



\emph{The groups.}  The isometry groups of the standard forms over $\R$, $\C$, $\HH$ are $O(n)$, $U(n)$, $Sp(n)$, and Theorem \ref{c:dark-tower-matter.constituents-field} assigns the fields. The script checks an $\F$-unitary mixing in every causal tower as a polynomial identity (C1) and exhibits a non-unitary mixing that changes the total momentum (C2).

\end{proof}

\phantomsection\label{c:note.dark-tower-matter.setting-ladder-pairs-complex-unit}
A \emph{ladder pair} is a pair $w, v$ of null vectors of the conformal space, $Q(w) = Q(v) = 0$, or of vectors dressed by a complex unit with $w^{2} = v^{2} = 0$, such that $wv + vw = 1$: the canonical anticommutation relations, one fermionic mode per pair \cite[Prop.~4.1]{ButlerT3}. Pairs built from real hyperbolic planes of $\hatV$ need no unit; their number is the real Witt index $\min(p,q) + 1$. The constituent's pair count is the complex Witt index $\lfloor (n+2)/2 \rfloor$ (T12's Lemma 2.12 in even dimension \cite[Lem.~2.12]{ButlerT12}; in odd dimension the same bound, with one generator left over).

\begin{proposition}[The ladder pairs of the causal towers]\label{c:dark-tower-matter.pairs}
\changedmark{BurntOrange}{changed since the deposit}
\begin{enumerate}
\item[(i)] In the three-dimensional towers the real and complex Witt indices of $\hatV$ are both $2$: two real ladder pairs, and no unit is needed (check D1).
\item[(ii)] In this tower, the mirror and the five-dimensional towers the real Witt index is $2$ and the complex one $3$: the third pair needs a complex unit. In this tower and the mirror the field is $\C$ and the unit is fixed up to sign (Remark \emph{Why U} [T14c, Rem 2.10]).
\end{enumerate}
\end{proposition}
\begin{proof}
In a causal tower $\min(p,q) = 1$, and $\hatV$ has signature $(p+1, q+1)$. Its Witt indices are
\[
\text{real: } \min(p+1, q+1) = 2, \qquad
\text{complex: } \Bigl\lfloor \frac{\dim\hatV}{2} \Bigr\rfloor =
\begin{cases} 2, & \dim\hatV = 5, \\ 3, & \dim\hatV = 6, 7. \end{cases}
\]
Check D1 builds the two real pairs of $\Cl(3,2)$ and $\Cl(2,3)$ on their modules and verifies the relations. In this tower and the mirror the field of (ii) is $\C$ by Theorem \ref{c:dark-tower-matter.constituents-field}.

\end{proof}

\begin{theorem}[The third pair's unit in five dimensions]\label{c:dark-tower-matter.third-pairs-unit-five-dimensions}
\changedmark{BurntOrange}{changed since the deposit}
Work on the sixteen-dimensional irreducible module of the conformal algebra, which is the constituent ideal $T$ as a module of the spin group, with spin commutant $\HH$ (Theorem \ref{c:dark-tower-matter.constituents-field}).
\begin{enumerate}
\item[(i)] ($r = 3$) The pseudoscalar $I_{7}$ of $\Cl(5,2)$ (here the pseudoscalar of the conformal algebra of the $(4,1)$ tower, not the unit $I_{7} = I_{6}e$ of $\Cl(3,4)$ that complexifies this tower's conformal algebra) lies in $\HH$, squares to $-1$ and commutes with every vector. With it the third pair $\tfrac{1}{2}(e_{a} \pm I_{7}e_{b})$ joins two real pairs to give three pairs obeying the canonical relations, and their bilinears $I_{7}(v_{i}w_{i} - \tfrac{1}{2})$, $v_{i}w_{k} - v_{k}w_{i}$, $I_{7}(v_{i}w_{k} + v_{k}w_{i})$ span a Lie algebra $\uu(3)$ whose centraliser in $\HH$ is $\mathrm{span}\{1, I_{7}\}$.
\item[(ii)] ($r = 3$) Every unit $J \in \HH$ anticommuting with $I_{7}$ anticommutes with every vector, so it acts as the grade involution: it fixes the real pairs up to sign and sends $w_{3}$ to $-\bar{w}_{3}$, exchanging the third pair's creator and annihilator; $\uu(3)$ and $J$ generate a Lie algebra of dimension $30$, not the commuting sum $\su(2) \oplus \uu(3)$ of dimension $12$.
\item[(iii)] ($r = 5$) No element of $\HH$ outside $\R$ commutes with all of $\HH$, and $\omega$ squares to $+1$: the algebra supplies no canonical unit. Every unit $u \in \HH$ commutes with the vectors and builds a third pair, $w_{3} = \tfrac{1}{2}(e_{a} + ue_{b})$, $v_{3} = -\tfrac{1}{2}(e_{a} - ue_{b})$ with $Q(e_{a}) = Q(e_{b}) = -1$; the resulting $\uu(3)_{u}$ has centraliser $\mathrm{span}\{1, u\}$ in $\HH$, and a second unit anticommuting with $u$ generates dimension $30$ with it. Any two units are conjugate in $\HH$: $(1 - u'u)\,u = u'(1 - u'u)$ for $u' \neq -u$, and $u' = -u$ is reached by conjugation with any unit anticommuting with $u$. The choices form the sphere $Sp(1)/U(1) = S^{2}$.
\end{enumerate}
\end{theorem}
\begin{proof}
\textit{(i) The unit $I_{7}$.}\enspace $I_{7}$ is central in $\Cl(5,2)$ (odd dimension) with $Q(I_{7}) = -1$ (the section ``Setting, imports and premises''). Central elements commute with the spin action, so $I_{7} \in \HH$, the spin commutant of Theorem \ref{c:dark-tower-matter.constituents-field}. Take $Q(e_{a}) = Q(e_{b}) = +1$ and
\[
w = \tfrac{1}{2}(e_{a} + I_{7}e_{b}), \qquad \bar{w} = \tfrac{1}{2}(e_{a} - I_{7}e_{b}).
\]
Since $I_{7}$ commutes with vectors and $e_{a}e_{b} = -e_{b}e_{a}$,
\begin{align*}
w^{2} &= \tfrac{1}{4}\bigl(Q(e_{a}) + I_{7}^{2}\,Q(e_{b})\bigr) = \tfrac{1}{4}(1 - 1) = 0, \\
w\bar{w} + \bar{w}w &= \tfrac{1}{2}\bigl(Q(e_{a}) - I_{7}^{2}\,Q(e_{b})\bigr) = \tfrac{1}{2}(1 + 1) = 1.
\end{align*}
The cross relations with the real pairs hold because distinct basis vectors anticommute and $I_{7}$ commutes with them (D3). Closure and the centraliser are computed (D4, D5).



\emph{(ii) A unit anticommuting with $I_{7}$.}  $J$ commutes with the even algebra (it lies in the spin commutant) and anticommutes with $I_{7}$. Every vector factors as $e = (eI_{7}^{-1})\,I_{7}$ with $eI_{7}^{-1}$ even, so
\[
JeJ^{-1} \;=\; (eI_{7}^{-1})\,\bigl(JI_{7}J^{-1}\bigr) \;=\; (eI_{7}^{-1})(-I_{7}) \;=\; -e \qquad \text{(D6)}.
\]
Hence $J(I_{7}e_{b})J^{-1} = I_{7}e_{b}$ and
\[
Jw_{3}J^{-1} \;=\; -\tfrac{1}{2}(e_{a} - I_{7}e_{b}) \;=\; -\bar{w}_{3},
\]
and the closure has dimension $30$ (D7).



\emph{(iii) The residue-five tower.}  The centre of $\HH$ is $\R$ and $Q(\omega) = +1$ (A4), so the algebra names no unit. The spin and Clifford commutants coincide on this module (D8), so every unit $u \in \HH$ commutes with the vectors. With $Q(e_{a}) = Q(e_{b}) = -1$ and $u^{2} = -1$,
\[
w_{3}v_{3} + v_{3}w_{3} \;=\; -\tfrac{1}{2}\bigl(Q(e_{a}) - u^{2}\,Q(e_{b})\bigr) \;=\; -\tfrac{1}{2}(-1 - 1) \;=\; 1 \qquad \text{(D9 to D11)}.
\]
For the conjugacy, with $u^{2} = u'^{2} = -1$,
\begin{align*}
(1 - u'u)\,u &= u - u'u^{2} = u + u', \\
u'\,(1 - u'u) &= u' - u'^{2}u = u' + u,
\end{align*}
so the two sides agree, and $1 - u'u = 0$ only at $u' = -u$ (D12).

\end{proof}

\section{The ladder pairs and the complex unit}

\begin{remark}[The two readings of the unit]\label{c:note.dark-tower-matter.two-readings-unit}
Theorem \ref{c:dark-tower-matter.third-pairs-unit-five-dimensions} is where the two five-dimensional towers part. At $r = 3$ the unit that builds the pairs is fixed by the algebra, and the quaternionic unit group of a single constituent, which preserves its momentum (Theorem \ref{c:dark-tower-matter.ladder}), does not preserve the pair structure. At $r = 5$ every choice of unit is equivalent, and each choice cuts the quaternionic group to the circle it generates. Under U3 both towers therefore build their pairs with one complex unit; they differ in whether the algebra names it. With an external unit, which commutes with the whole quaternionic group, the group survives in both towers alike (Remark \ref{c:note.dark-tower-matter.branches-break-u-what-they-give}).
\end{remark}

\begin{theorem}[The gauge algebras of the odd towers]\label{c:dark-tower-matter.gauge-algebras}
\changedmark{BurntOrange}{changed since the deposit}
Under H2 and Premise \ref{c:dark-tower-matter.u-transposition-rule}:
\begin{enumerate}
\item[(i)] (Three-dimensional towers, $r = 1, 7$.) Mass begins at two real constituents and each carries two real pairs, so the minimal massive hierarchy is $\R^{2} \otimes \R^{2}$. Its automorphisms $A \otimes B$ with $A, B \in O(2)$ form $(O(2) \times O(2))/\mathbb{Z}_{2}$, the kernel being $\{(1,1), (-1,-1)\}$; the gauge algebra is $\so(2) \oplus \so(2)$, abelian, of dimension $2$, and the identity component of the group is $U(1)_{+} \times U(1)_{-}$.
\item[(ii)] (Five-dimensional towers, $r = 3, 5$.) Mass begins at one constituent and its three pairs are complex with the unit of Theorem \ref{c:dark-tower-matter.third-pairs-unit-five-dimensions}, so the hierarchy is $\C^{1} \otimes \C^{3}$. Its automorphisms form $(U(1) \times U(3))/U(1) \cong U(3)$; the gauge algebra is $\uu(3) = \su(3) \oplus \uu(1)$, of dimension $9$.
\end{enumerate}
\end{theorem}
\begin{proof}
\textit{(i) The two factors.}\enspace The constituent count is two (Theorem \ref{c:dark-tower-matter.where-mass-begins}(i)) and the field is $\R$ (Theorem \ref{c:dark-tower-matter.constituents-field}), so by U2 (\ref{c:dark-tower-matter.u-transposition-rule}) the slot group is $O(2)$ (Theorem \ref{c:dark-tower-matter.ladder}). The pairs are real (Proposition \ref{c:dark-tower-matter.pairs}(i)), and by U2 their group is the orthogonal group of the real pair space, again $O(2)$.

The norm clause of U2 matters here. Without it, the pair mixings preserving the canonical relations alone would be $GL(2,\R)$, acting by $B$ on the $w_{i}$ and by $B^{-T}$ on the $v_{i}$. That is the mechanism of Remark \ref{c:note.dark-tower-matter.kleinian-algebra}, where the sterile tower has no time direction and so no adjoint.



\emph{(i) The kernel.}  An automorphism of $\R^{2} \otimes \R^{2}$ respecting both factors is $A \otimes B$, and
\[
A \otimes B = 1 \quad\Longleftrightarrow\quad A = \lambda 1, \;\; B = \lambda^{-1}1,
\]
with $\lambda$ real and orthogonal, so $\lambda = \pm 1$. The kernel $\{(1,1), (-1,-1)\}$ is discrete, no generator is lost, and the algebra is $\so(2) \oplus \so(2)$ (the gauge algebra by \ref{c:dark-tower-matter.h2-identification}; checks E1, E2).



\emph{(i) The charges.}  Write $s$ and $p$ for the weights of the slot and pair rotations. The element $(-1, -1)$ acts as $(-1)^{s + p}$, so the group's representations are those with $s + p$ even, on which
\[
q_{\pm} = \frac{s \pm p}{2}
\]
are the integral charges of $U(1)_{\pm}$ (F2).



\emph{(ii) The slot.}  The constituent count is one (Theorem \ref{c:dark-tower-matter.where-mass-begins}(ii)). The slot group of one constituent is $Sp(1)$ (Theorem \ref{c:dark-tower-matter.ladder}), acting on the spinor part through H1 (\ref{c:dark-tower-matter.h1-momentum})'s identification of the two fields. The pair unit of Theorem \ref{c:dark-tower-matter.third-pairs-unit-five-dimensions} lies in that field, and only the circle it generates commutes with the pair algebra, so the slot is one complex line with group $U(1)$.



\emph{(ii) The pairs and the kernel.}  By U2 the three complex pairs carry $U(3)$. The action $(\lambda, B) \mapsto \lambda B$ has the common scalar as its kernel, as in T12's proof of its Theorem 4.1, so the group is
\[
\frac{U(1) \times U(3)}{U(1)} \;\cong\; U(3)
\]
(E3; the control $\C^{2} \otimes \C^{3} \mapsto 12$ is E4).

\end{proof}

\phantomsection\label{c:note.dark-tower-matter.setting-gauge-algebras}
In the mirror the algebra and global form are this tower's, $\su(2) \oplus \su(3) \oplus \uu(1)$ and $S(U(2) \times U(3))$: that is T14's copy theorem (\ref{c:dark-matter.mirror-tower-copy-one}), which Theorem \ref{c:dark-tower-matter.mirror-algebra} below realises as an explicit isomorphism of the even algebras.

\section{The gauge algebras}

\begin{remark}[Branches that break U, and what they give]\label{c:note.dark-tower-matter.branches-break-u-what-they-give}
\changedmark{gray}{reworded}
Two transpositions reproduce this tower and break a clause of Premise \ref{c:dark-tower-matter.u-transposition-rule}. (a) Counting the $r = 7$ constituent in the full Clifford module, where the central $\omega$ joins two real constituents into one complex element that is already massive (Theorem \ref{c:dark-tower-matter.where-mass-begins}(i) for the two real parts), gives the hierarchy $\C^{1} \otimes \C^{2}$ and $\uu(2)$, with content a doublet of $Y = \tfrac{1}{2}$, a singlet of $Y = -1$ and a neutral singlet; it breaks U1, since $\omega$ is odd and does not act on the spin-irreducible ideal (Proposition \ref{c:dark-tower-matter.residue-seven-unit}(ii)). (b) Building the five-dimensional third pair with an external unit, one adjoined from outside the real algebra, which commutes with the quaternionic unit group, keeps that group as the slot group of the one constituent: the hierarchy is a quaternionic line tensor $\C^{3}$, with automorphisms $\su(2) \oplus \su(3) \oplus \uu(1)$ and global form $S(U(2) \times U(3))$, and vector-like content $(2,3)_{1/6}$, $(2,1)_{-1/2}$. It keeps U1 and U2 and breaks U3; it is T12's own practice of working in the complexification, carried to a tower whose field is $\HH$; and it is open at $r = 3$ and $r = 5$ on the same terms, so it does not separate them. Its $\su(2)$ acts on one constituent's quaternionic multiplicity, not on two constituents. Their arithmetic is in the supplementary scripts (\file{D3C}, \file{D5Q}; the section ``Machine verification''). Neither is a result of this paper.
\end{remark}

\phantomsection\label{c:note.dark-tower-matter.setting-content}
T12's lesson is that fermions are generators of the composite, not its states \cite[Rem.~5.2, Prop.~5.3]{ButlerT12}. The generators of a hierarchy $\F^{m} \otimes \F^{k}$ are its single slot-pair cells; Hypothesis (S) adds the constituent-level states.

\begin{proposition}[The three-dimensional towers' content]\label{c:dark-tower-matter.three-dimensional-content}
\changedmark{BurntOrange}{changed since the deposit}
In the towers $r = 1, 7$:
\begin{enumerate}
\item[(i)] The generators are the four cells of $\R^{2} \otimes \R^{2}$, of weights $(s, p) \in \{\pm 1\}^{2}$; as complex two-component fermions they are $A$, with $(q_{+}, q_{-}) = (1, 0)$, and $B$, with $(0, 1)$, and their conjugates. Each is charged under one factor only: the tower is two copies of single-flavour electrodynamics in $2+1$ dimensions, exchanged by the pair reflection $p \mapsto -p$.
\item[(ii)] Every weight of the exterior algebra of the cells has $s + p$ even; no state of the tensor state space has weight $(\pm 1, 0)$ or $(0, \pm 1)$.
\item[(iii)] Under (S) the constituent-level states have weight $(\pm 1, 0)$, off the lattice $s + p$ even: (S) and the global form of Theorem \ref{c:dark-tower-matter.gauge-algebras}(i) are incompatible. Of the directions the transport adds, the rung directions carry $(\pm 1, 0)$ and the pair directions $(0, \pm 1)$, and both are excluded; the one neutral direction (the leftover generator of the odd conformal algebra) is allowed, and its value is a bare mass. The tower has no charged Higgs analogue.
\item[(iv)] There is no chirality: the odd-dimensional modules carry no Lorentz-invariant grading \cite[Prop.~3.9]{ButlerT14}. Gauge-invariant Dirac masses exist for $A$ and $B$ without any condensate, and there is no baryon-like composite: every gauge-neutral state has zero net number of each cell.
\end{enumerate}
\end{proposition}
\begin{proof}
\textit{(i)}\enspace The cells are single-occupation states. The charges $q_{\pm} = (s \pm p)/2$ of the $U(1)_{+} \times U(1)_{-}$ of Theorem \ref{c:dark-tower-matter.gauge-algebras}(i) send
\[
(s, p) = (1, 1) \;\mapsto\; (q_{+}, q_{-}) = (1, 0), \qquad (1, -1) \;\mapsto\; (0, 1) \qquad \text{(F2)},
\]
and the pair reflection $p \mapsto -p$ exchanges them.



\emph{(ii)}  A product of cells adds weights, and each cell has $s + p \in \{-2, 0, 2\}$ (F1).



\emph{(iii)}  A constituent-level state occupies one slot mode and no pair mode, so its weight is $(\pm 1, 0)$, off the lattice $s + p$ even (Remark \ref{c:note.dark-tower-matter.why-s-fails-three-dimensions-not}). The weights of the transported directions are computed in the supplementary script \file{D3R} (checks H5 to H9): rung directions $(\pm 1, 0)$, pair directions $(0, \pm 1)$, the leftover generator $(0,0)$.



\emph{(iv)}  This is the odd case of T14's Proposition \ref{c:dark-matter.two-families-are-two-chiralities}. A mass term $\bar{A}A$ has weight zero, and a neutral product of cells pairs each cell with its conjugate.

\end{proof}

\begin{remark}[The neutral row]\label{c:note.dark-tower-matter.neutral-row}
\changedmark{BurntOrange}{changed since the deposit}
In a hierarchy $\C^{1} \otimes \C^{k}$ whose central generator acts by $k$ on the slot and $-1$ on each pair, the state with the slot and all $k$ pairs occupied has weight $k - k = 0$. The constituent-level generator of a one-constituent hierarchy is therefore neutral: the trace-free central generator of $\uu(1) \oplus \uu(k)$ on $\C^{1} \oplus \C^{k}$ (numbers and the module's complex structure are read in $\mathrm{span}\{1, I_{4}\}$, the module read on the tower's floors, where $I_{4}$ commutes with the tower vectors) is $(k\,I_{4}, -I_{4}\,1_{k})$ up to scale, and weights add. In the five-dimensional towers the pairs themselves are built with the unit $u$ of the hierarchy's scalars $\mathrm{span}\{1, u\}$ ($u = I_{7}$, the pseudoscalar of $\Cl(5,2)$, at $r = 3$; a chosen unit of $\HH$ at $r = 5$; Theorem \ref{c:dark-tower-matter.third-pairs-unit-five-dimensions}); here $I_{4}$ is the unit of the multiplicity spaces only, and the two are not identified. Arithmetic.
\end{remark}

\begin{proposition}[The five-dimensional towers' content]\label{c:dark-tower-matter.five-dimensional-content}
\changedmark{BurntOrange}{changed since the deposit}
In the towers $r = 3, 5$, under (S) the state space $\Lambda^{\bullet}(\C^{1} \oplus \C^{3})$ has sixteen states. With $Y = s - k/3$ ($s$ the slot occupation, $k$ the number of occupied pairs), the complement $(s,k) \mapsto (1-s, 3-k)$ negates $Y$ and conjugates colour, so the sixteen are eight species-states and their conjugates, four vector-like Dirac species:
\begin{center}\small
\begin{tabular}{@{}lccl@{}}
\toprule
species & colour & $Y$ & states $(s,k)$ \\
\midrule
$U$ & $3$ & $2/3$ & $(1,1)$, conjugate $(0,2)$ \\
$D$ & $3$ & $-1/3$ & $(0,1)$, conjugate $(1,2)$ \\
$E$ & $1$ & $-1$ & $(0,3)$, conjugate $(1,0)$ \\
$N$ & $1$ & $0$ & $(0,0)$, $(1,3)$ \\
\bottomrule
\end{tabular}
\end{center}
These are the electric charges of one of our generations with the weak doublet removed. Through the section $B \mapsto ((\det B)^{-1}, B)$ the centre of $U(3)$ acts on $(s,k)$ by $\lambda^{k - 3s} = \lambda^{-3Y}$, so $3Y$ is an integer with $3Y + k \equiv 0 \pmod 3$, $k$ being the triality: the lattice of $U(3)$, satisfied automatically; the neutral species $N$ is the empty state and the full row of Remark \ref{lem:neutralrow}. Without (S) the content is one Dirac colour triplet of $Y = 2/3$. There is no chirality and every species has a gauge-invariant bare mass.
\end{proposition}
\begin{proof}
The state space under (S) (\ref{c:dark-tower-matter.hypothesis-s-slot-statistics}) has
\[
\dim \Lambda^{\bullet}\C^{4} \;=\; \sum_{s=0}^{1}\sum_{k=0}^{3}\binom{3}{k} \;=\; 2 \cdot 8 \;=\; 16
\]
states, the $\binom{3}{k}$ states at pair number $k$ forming the $1, 3, \bar{3}, 1$ of the $SU(3)$ of Theorem \ref{c:dark-tower-matter.gauge-algebras}(ii). The complement negates $Y$:
\[
Y(1 - s, 3 - k) \;=\; (1 - s) - \frac{3 - k}{3} \;=\; -\Bigl(s - \frac{k}{3}\Bigr) \;=\; -Y(s, k),
\]
and the remaining weights are arithmetic (F3, F4). The masses are as in Proposition \ref{c:dark-tower-matter.three-dimensional-content}(iv).

\end{proof}

\section{Content}

\begin{remark}[Why (S) fails in three dimensions, not five]\label{c:note.dark-tower-matter.why-s-fails-three-dimensions-not}
In both cases the slot-level states are charged under the kernel of the action on the hierarchy. In five dimensions the kernel is the common scalar $U(1)$, and the subgroup $S(U(1) \times U(3)) = \{((\det B)^{-1}, B)\}$ maps isomorphically onto the quotient $U(3)$: the exterior algebra of slot and pair modes is a representation of that section, as T12 uses $S(U(2) \times U(3))$ in four dimensions. In three dimensions the kernel $\{\pm(1,1)\}$ is discrete inside the torus $SO(2) \times SO(2)$, and no subgroup maps isomorphically onto the quotient, so the slot-level states are not representations of the global form. Under (S) the five-dimensional state space has a tenth abelian automorphism, the slot number, beyond the nine of $\uu(3)$; this paper takes it as a global symmetry, not gauged (a choice: T12 counts the corresponding extra generator in four dimensions and finds it anomaly-free, broken by the right-handed neutrino's Majorana mass \cite[Prop.~5.8]{ButlerT12}); here it is broken only by a mass term for $N$ off the pair unit's axis.
\end{remark}

\begin{proposition}[Chern-Simons levels of the odd towers' content]\label{c:dark-tower-matter.anomalies}
\changedmark{BurntOrange}{changed since the deposit}
\begin{enumerate}
\item[(i)] In the three-dimensional towers each $U(1)_{\pm}$ couples to exactly one unit-charge two-component fermion, which shifts its Chern-Simons level by $\tfrac{1}{2}\mathrm{sign}(m)$ \cite{Redlich1984}; integrality of the effective level forces the bare levels into $\mathbb{Z} + \tfrac{1}{2}$. Parity is therefore broken by the gauge sector.
\item[(ii)] In the five-dimensional towers, in units $q' = 3Y$, the content of Proposition \ref{c:dark-tower-matter.five-dimensional-content} has $\sum q'^{3} = -6$ and $\sum q' = 0$ (with multiplicities), and the induced abelian cubic Chern-Simons coefficient, $\tfrac{1}{2}\sum \mathrm{sign}(m)\,q'^{3}$, is integral for every choice of mass signs \cite{IMS1997}. The $SU(3)$ and mixed levels, which on $U(3)$ are tied to the abelian one, are not analysed here.
\end{enumerate}
\end{proposition}
\begin{proof}
\textit{(i)}\enspace With bare level $k$ and one unit-charge fermion of mass $m$ (Proposition \ref{c:dark-tower-matter.three-dimensional-content}(i)), the effective level is, by Fact \ref{c:rule2.imported-parity-anomaly},
\[
k_{\mathrm{eff}} \;=\; k + \tfrac{1}{2}\,\mathrm{sign}(m) \;\in\; \mathbb{Z} \quad\Longrightarrow\quad k \in \mathbb{Z} + \tfrac{1}{2} \qquad \text{(G1, G2)},
\]
and no half-integer is zero, so a parity-odd Chern-Simons term is always present.



\emph{(ii)}  In units $q' = 3Y$ the species of Proposition \ref{c:dark-tower-matter.five-dimensional-content} carry $q' = 2, -1, -3, 0$ for $U, D, E, N$, with colour multiplicity $3$ for $U$ and $D$. So
\begin{align*}
\sum q'^{3} &= 3(2)^{3} + 3(-1)^{3} + (-3)^{3} + 0 = 24 - 3 - 27 = -6, \\
\sum q' &= 3(2) + 3(-1) + (-3) + 0 = 0 \qquad \text{(G3)}.
\end{align*}
In the induced coefficient $\tfrac{1}{2}\sum \mathrm{sign}(m)\,q'^{3}$ (Fact \ref{c:rule2.imported-five-dimensional-chern-simons-shift}) the terms of $U$ and $N$ are the integers $\pm 12$ and $0$, and those of $D$ and $E$ are $\pm\tfrac{3}{2}$ and $\pm\tfrac{27}{2}$, whose sum is an integer for each choice of signs.

\end{proof}

\phantomsection\label{c:note.dark-tower-matter.setting-paired-towers}
T14's indifference premise P2 equates the data of the two members of each reversal pair, $\{2, 6\}$, $\{1, 7\}$ and $\{3, 5\}$ \cite[Prem.~2.2]{ButlerT14}. This section compares the members of the odd pairs as algebras.

\begin{proposition}[The central unit of the residue-seven tower]\label{c:dark-tower-matter.residue-seven-unit}
\changedmark{gray}{reworded}
\begin{enumerate}
\item[(i)] Along the residue-seven tower ($\Cl(2,3)$, $\Cl(3,4)$, $\Cl(4,5)$, \ldots) the pseudoscalar is central and squares to $-1$ at every rung; along the residue-one tower ($\Cl(3,2)$, $\Cl(4,3)$, \ldots) it is central and squares to $+1$.
\item[(ii)] At $r = 7$ the pseudoscalar $\omega$ is odd: it carries the even ideal into the odd one, so it does not act on the constituent ideal, and the grade involution reverses it.
\item[(iii)] On the Clifford module $\Cl(\hatV)\hf = T \oplus \omega T$, left multiplication by $c + s\omega$ ($c^{2} + s^{2} = 1$) commutes with the whole Clifford action and carries $X + \omega Y$ to $(cX - sY) + \omega(sX + cY)$: on a pair of real constituents it is the slot rotation of Theorem \ref{c:dark-tower-matter.gauge-algebras}(i), here supplied by the algebra itself. Every cell carries charge $\pm 1$ under it. If the tower's action is built from operations that commute with the Clifford action, as T12's are, a Majorana-type mass (charge $2$) is then forbidden in the $r = 7$ tower whether or not the rotation is gauged; this last clause is at assembly level.
\end{enumerate}
\end{proposition}
\begin{proof}
\textit{(i)}\enspace By the formula of the section ``Setting, imports and premises'' (\ref{c:primer.sign}), $Q(\omega) = (-1)^{n(n-1)/2}(-1)^{q}$. Each rung adds $(1,1)$ to the signature, so $n$ rises by two and $q$ by one, and the first exponent changes by
\[
\frac{(n+2)(n+1)}{2} - \frac{n(n-1)}{2} \;=\; 2n + 1,
\]
which is odd. The two sign changes cancel, so $Q(\omega)$ is constant along each tower and equal to its value at the first rung:
\[
Q(\omega_{2,3}) = (-1)^{10}(-1)^{3} = -1, \qquad Q(\omega_{3,2}) = (-1)^{10}(-1)^{2} = +1.
\]
Every rung has odd dimension, so $\omega$ is central (H1, H2; \ref{c:std-algebra.pseudoscalar-commutation}).



\emph{(ii)}  $\omega$ is a product of five vectors (H3), so it is odd, left multiplication by it exchanges the even and odd parts, and $\hat{\omega} = (-1)^{5}\omega = -\omega$ (\ref{c:std-algebra.involution-properties}).



\emph{(iii)}  $\omega$ is central with $\omega^{2} = -1$ (H4), so $c + s\omega$ commutes with the Clifford action, and for $X, Y \in T$
\[
(c + s\omega)(X + \omega Y) \;=\; cX + s\omega^{2}Y + \omega(sX + cY) \;=\; (cX - sY) + \omega(sX + cY).
\]

\end{proof}

\begin{corollary}[Paired towers hold the same kind of matter]\label{c:dark-tower-matter.same-kind-matter}
\changedmark{BurntOrange}{changed since the deposit}
Under H2 and Premise \ref{c:dark-tower-matter.u-transposition-rule} the two towers of each odd dimension have the same constituent field, the same mass onset, the same gauge algebra and the same content. They differ only as follows.
\begin{center}\small
\begin{tabular}{@{}p{3.6cm}p{4.9cm}p{4.9cm}@{}}
\toprule
 & $r = 1$ & $r = 7$ \\
\midrule
slot rotation & a rotation of two modules of the same type & the algebra's central unit, exact at every rung \\
slot reflection & internal & exists only fused with orientation reversal \\
module types & two ($\omega = \pm 1$) & one \\
number-violating masses & forbidden once the slot rotation is gauged (H2) & forbidden by the algebra's unit, gauged or not \\
\midrule
 & $r = 3$ & $r = 5$ \\
\midrule
the pair unit & $I_{7}$, fixed by the algebra & a point of $S^{2}$, not fixed \\
reversion-type invariant form & one, Hermitian for $I_{7}$ & none \\
module types & one & two \\
\bottomrule
\end{tabular}
\end{center}
\end{corollary}
\begin{proof}
The shared rows are Theorems \ref{c:dark-tower-matter.constituents-field}, \ref{c:dark-tower-matter.where-mass-begins}, \ref{c:dark-tower-matter.gauge-algebras} and Propositions \ref{c:dark-tower-matter.three-dimensional-content}, \ref{c:dark-tower-matter.five-dimensional-content}. The differing rows come from Proposition \ref{c:dark-tower-matter.residue-seven-unit} and Theorem \ref{c:dark-tower-matter.third-pairs-unit-five-dimensions}, except two that are computed in supplementary scripts: the slot-reflection row in \file{D3C} (Part I), and the invariant-form row in the referee script \file{CHECK\_ALG/check\_r3\_r5.py}.

The argument holds given \ref{c:dark-tower-matter.u-transposition-rule}, which fixes the constituent field and the slot and pair groups, and \ref{c:dark-tower-matter.h2-identification}, which identifies the gauge algebra with the automorphism algebra of the minimal massive hierarchy.

\end{proof}

\section{The paired towers}

\begin{remark}[The $r = 5$ unit: datum or field]\label{c:note.dark-tower-matter.r-5-unit-datum-or-field}
If the unit $u$ of Theorem \ref{c:dark-tower-matter.third-pairs-unit-five-dimensions}(iii) is a constant of the tower, every choice gives an isomorphic tower and nothing observable distinguishes it from $r = 3$; this is how this tower treats its own unit, which the algebra fixes up to sign. If $u$ varies over spacetime, the tower carries a sigma model on $S^{2}$ with two massless five-dimensional scalars, which moves the radiation bound of Computation \ref{c:dark-tower-matter.five-dimensional-towers}(iii) from $0.085$ to $0.084$. The series supplies no kinetic term for $u$; the first reading is this tower's.
\end{remark}

\begin{theorem}[The mirror's even algebra as a real isomorphism]\label{c:dark-tower-matter.mirror-algebra}
\changedmark{BurntOrange}{changed since the deposit}
\begin{enumerate}
\item[(i)] The linear map $e_{K} \mapsto (-1)^{|K|/2}e_{K}$ on even basis blades is an algebra isomorphism from the even algebra of a signature onto the even algebra of the reversed signature (every $Q(e_{a})$ negated); in particular $\Cl^{0}(3,1) \cong \Cl^{0}(1,3)$ and $\Cl^{0}(4,2) \cong \Cl^{0}(2,4)$.
\item[(ii)] It carries the mirror's pseudoscalar $I_{4}'$ to $+I_{4}$ and its conformal pseudoscalar $I_{6}'$ to $-I_{6}$.
\item[(iii)] The full algebras differ: $\Cl(1,3)$ has quaternionic irreducible modules and $\Cl(3,1)$ real ones.
\end{enumerate}
\end{theorem}
\begin{proof}
\textit{(i) The product rule.}\enspace Write $Q$ for the squares of the first signature and $Q' = -Q$ for those of the second. For even blades $e_{K}$, $e_{L}$,
\[
e_{K}e_{L} \;=\; \sigma \prod_{c \in K \cap L}Q(e_{c})\;e_{K \triangle L},
\]
with $\sigma$ the reordering sign, which does not depend on the squares; the same holds in the second algebra with $Q'$ in place of $Q$.



\emph{(i) The map respects it.}  Let $\varphi(e_{K}) = (-1)^{|K|/2}e_{K}$, and use $|K \triangle L| = |K| + |L| - 2|K \cap L|$ and $\prod_{c}Q'(e_{c}) = (-1)^{|K \cap L|}\prod_{c}Q(e_{c})$. Then
\begin{align*}
\varphi(e_{K}e_{L}) &= \sigma\prod_{c}Q(e_{c})\,(-1)^{|K \triangle L|/2}\,e_{K \triangle L}
= (-1)^{\frac{|K|+|L|}{2}}(-1)^{|K \cap L|}\,\sigma\prod_{c}Q(e_{c})\,e_{K \triangle L}, \\
\varphi(e_{K})\,\varphi(e_{L}) &= (-1)^{\frac{|K|+|L|}{2}}\,\sigma\prod_{c}Q'(e_{c})\,e_{K \triangle L}
= (-1)^{\frac{|K|+|L|}{2}}(-1)^{|K \cap L|}\,\sigma\prod_{c}Q(e_{c})\,e_{K \triangle L}.
\end{align*}
The two agree. The script checks every pair of basis blades in both cases (I1).



\emph{(ii)}  On the pseudoscalars,
\[
\varphi(I_{4}') = (-1)^{4/2}\,I_{4} = +I_{4}, \qquad \varphi(I_{6}') = (-1)^{6/2}\,I_{6} = -I_{6} \qquad \text{(I2)}.
\]



\emph{(iii)}  The commutants are $\HH$ and $\R$ (I3; Fact \ref{c:dark-tower-matter.imported-clifford-structure}).

\end{proof}

\phantomsection\label{c:note.dark-tower-matter.setting-mirror}
So the mirror's constituents are complex, as ours are (Theorem \ref{c:dark-tower-matter.constituents-field}): its real (Majorana) ground spinors do not make them real. The isomorphism sends the mirror's field unit $I_{4}'$ to ours and its conformal unit $I_{6}'$ to minus ours. If each tower's constituent unit is taken to be its own oriented conformal pseudoscalar, the mirror's constituent unit and field unit therefore have the opposite relative sign to ours, and the mirror is the parity image of this tower. That reading is interpretive and depends on the identification of the transported plane, which with one sign reversed would flip the sign of $I_{6}'$'s image. Nothing gravitational distinguishes the two.

\section{The mirror}

\begin{proposition}[The mirror's temperature and premise P2]\label{c:dark-tower-matter.mirrors-temperature-p2}
\changedmark{BurntOrange}{changed since the deposit}
T14 applies P2 along the reversal orbit to equate the ceiling coupling and the asymmetry per mode, both data of the descent \cite[Thm.~4.3(iv), (v)]{ButlerT14}. The energy the descent deposits per unit charge is a datum of the same kind, and P2 declares which of $2$, $6$ is visible to be indexical. Applied with the same reach, P2 gives $\xi = T'/T = 1$ and $\Delta\Neff = \tfrac{8}{7}(\tfrac{11}{4})^{4/3} + 3.044 = 7.45$ at recombination, excluded by a factor of about $25$. T14's mirror, cold from the gravitational floor upward \cite[Prop.~4.6, Prop.~A.2]{ButlerT14}, is therefore an exception to P2 that T14 neither states nor costs. It is an exception at P2's version-1 reach, the reach T14 used; P2 as scoped (version 2) leaves $\xi$ free (Theorem \emph{The tower charge is a sector label; the temperature is not}), and the cold mirror then needs a non-equilibrium descent (Proposition \emph{A cold mirror needs a non-equilibrium descent}). Assembly level.
\end{proposition}
\begin{proof}
\textit{The reach of P2.}\enspace The equalities of T14's Theorem \ref{c:dark-matter.mirror-tower-copy-one}(iv), (v), as T14 states them (the theorem's version 8; at its current version (iv) is the common-ceiling-coupling premise and (v) reads P2 on the stored charge), rest on P2's clause that a difference between two continuous data of one reversal orbit is an uncaused distinction. T14's Proposition A.2 (\ref{c:dark-matter.seeding-gravity-floor-dark-towers-temperature}) calls what the descent deposited ``a datum, not a computation'', and T14's Proposition 4.6 (\ref{c:dark-matter.stored-asymmetry-charge-not-number}) carries the mirror's temperature as such a datum, from the floor up to the $\Neff$ cap.



\emph{The number.}  At $\xi = 1$ the mirror is a full Standard-Model copy at the photon temperature. Its energy density is that of a photon, worth $\tfrac{8}{7}(\tfrac{11}{4})^{4/3}$ neutrino species, plus three neutrinos:
\[
\Delta\Neff \;=\; \underbrace{\tfrac{8}{7}\bigl(\tfrac{11}{4}\bigr)^{4/3}}_{\text{mirror photon}} \;+\; \underbrace{3.044}_{\text{mirror neutrinos}} \;=\; 4.40 + 3.044 \;=\; 7.45 \qquad \text{(K2)}.
\]

\end{proof}

\begin{theorem}[The mirror's composition across its temperature range]\label{c:dark-tower-matter.mirrors-composition}
\changedmark{BurntOrange}{changed since the deposit}
Decimal. Supplementary scripts: the mirror's (folder \file{MIR}, checks D1, D2 and E0 to E7) and the cosmology check's (folder \file{CHECK\_COS}). T14 carries $\xi$ from the gravitational floor to the $\Neff$ cap ($0.448$), with a hot-deposit estimate near $4 \times 10^{-3}$. At the top of the ceiling range the floor was $2.8$ to $3.2 \times 10^{-4}$ with the production coefficients first imported, and the computation below was first made there; with the derived coefficients and the mean energy per creation event it is $6.6 \times 10^{-4}$ (Proposition \ref{c:dark-matter.stored-asymmetry-charge-not-number}(iii)), and the two scripts were run again on 18 November 2026 with that value added to their grids.
\begin{enumerate}
\item[(i)] At $\xi = 2.8$ to $3.2 \times 10^{-4}$, the floor when this was computed and now the floor for a reheating below the top of the ceiling range (the floor goes as $T_{\mathrm{desc}}^{3/4}$), the mirror is charge-dominated and degenerate from the moment it is seeded: the Fermi energy its deposited baryon number must carry exceeds the thermal deposit there by $49$ to $84$, and at its weak freeze-out the quark chemical potential is $18$ to $23$ times its temperature. Beta equilibrium persists into the hadron phase, and for a copied lepton number $L \leq 0$ (our tower's sphaleron value $L/B = -51/28$, or zero) the mirror ends $96$ to $100$ percent neutrons. At $\xi = 6.6 \times 10^{-4}$, the present floor at the top of the ceiling range, the same scripts give a mirror that is still neutron-rich for $L \leq 0$ but only moderately degenerate: the Fermi energy exceeds the thermal deposit by $2.7$, the quark chemical potential is $3.5$ times the temperature at seeding and $3.7$ times at $T' \approx 0.4$ GeV (the quantity that is $18$ to $23$ at the lower values), beta equilibrium still lasts into the hadron phase, and the mirror ends $96$ percent neutrons for $L/B = -51/28$ and $83$ percent for $L = 0$. So ``charge-dominated and degenerate from the moment it is seeded'' holds at $\xi = 2.8$ to $3.2 \times 10^{-4}$ and holds at the present floor only in this weaker form: the charge still dominates, by a factor near three and not by fifty or more. At either value the endpoint is cold, dense, neutron-rich matter; the composition after decay and capture needs a dense-matter network and is not computed here.
\item[(ii)] Between $\xi \approx 3 \times 10^{-3}$ and $2 \times 10^{-2}$, an interval containing T14's hot-deposit estimate, the mirror is helium ($Y' = 0.94$ to $1.00$).
\item[(iii)] Above that interval the helium fraction falls ($Y' = 0.85$ at $\xi = 0.1$) and towards the cap the mirror is hydrogen-rich ($Y' = 0.58$ at $\xi = 0.448$).
\end{enumerate}
T14 does not print a mirror $\eta'$ at freeze-out; its floor value ($19$ to $26$ as printed, $2.1$ with the derived coefficients and the mean energy per creation event) is the late-time ratio, and the thermal up-down ratio T14 assumes there is inconsistent with its own $\eta' > 1$.
\end{theorem}
\begin{proof}
Computed, as the source states, in the supplementary scripts \file{MIR} (checks D1, D2 and E0 to E7) and \file{CHECK\_COS}, across the range of the temperature ratio that \ref{c:dark-matter.stored-asymmetry-charge-not-number} and \ref{c:dark-matter.seeding-gravity-floor-dark-towers-temperature} supply, for the mirror that \ref{c:dark-matter.mirror-tower-copy-one} makes a copy of this tower. The source gives no further proof text. Clause (i) was first computed at $\xi = 2.8$ to $3.2 \times 10^{-4}$, the floor of \ref{c:dark-matter.stored-asymmetry-charge-not-number} at the production coefficients first imported; the floor there is now $6.6 \times 10^{-4}$ (derived coefficients, mean energy per creation event). The statement at the present floor is from a second run of the same two scripts on 18 November 2026, each changed in one line, the grid of $\xi$, to which $5.2 \times 10^{-4}$, $6.6 \times 10^{-4}$ and $8.1 \times 10^{-4}$ were added; the mirror's script passes its $24$ checks as before and reproduces the figures at $2.8$ and $3.2 \times 10^{-4}$. At $\xi = 6.6 \times 10^{-4}$ the cosmology check gives the ratio of Fermi energy to thermal deposit as $2.7$ ($84$ and $49$ at $2.8$ and $3.2 \times 10^{-4}$) and the quark chemical potential over the temperature as $3.5$ at seeding and $3.7$ at $T' \approx 0.4$ GeV ($23.4$ and $17.8$ there). The mirror's script, with its degeneracy-aware weak freeze-out in the hadron phase, gives freeze-out at $T' = 18$ MeV with $n/p = 22$ and neutron fraction $0.957$ for $L/B = -51/28$, and at $T' = 111$ MeV with $n/p = 5.0$ and neutron fraction $0.832$ for $L = 0$ (at $2.8$ and $3.2 \times 10^{-4}$: $1.000$, $0.999$ and $0.965$, $0.954$); for a positive lepton number $L_{e}/B = 0.875$, outside the clause, it gives $0.32$. Its thermal-rate table, the paper's own method, puts the late-time ratio $\eta'$ at $2.12$ and the quark chemical potential at freeze-out at $1.8$ times the temperature, where it has $10$ and $8.6$ at the lower values. On the helium branch, which is not the branch the mirror takes there for $L \leq 0$, the same run puts the deuterium bottleneck at $T' = 0.185$ MeV with $\eta' = 2.13$ and a baryon density of $4.3 \times 10^{29}$ per cm$^{3}$, and the burning reheats the mirror's plasma by $1.86$ ($3.23$ and $2.96$, with $1.3 \times 10^{31}$ and $7.6 \times 10^{30}$ per cm$^{3}$, at $2.8$ and $3.2 \times 10^{-4}$). For comparison the run gives, at $\xi = 5.2 \times 10^{-4}$, the ratio $7.1$, neutron fractions $0.984$ and $0.887$ and a reheat of $2.16$, and at $8.1 \times 10^{-4}$, the ratio $1.2$, neutron fractions $0.916$ and $0.771$ and a reheat of $1.64$; the ratio of Fermi energy to thermal deposit passes one at $8.5 \times 10^{-4}$, the crossover of the cited Proposition. These are decimal results of scripts that are not in the store: no check of them is added, and the composition after decay and capture is not computed at any of these values.

\end{proof}

\begin{conjecture}[The mirror against the astrophysical record]\label{c:dark-tower-matter.mirror-against-record}
\changedmark{BurntOrange}{changed since the deposit}
On every branch the mirror recombines early (at a visible temperature of $0.04$ to $3$ keV) and is cold dark matter for the microwave background and large-scale structure, and in galaxies it is a collisional gas that cools by bremsstrahlung in a few gigayears, settles partly and forms stars. Against the Gaia thin-disk bounds \cite{Widmark2021} and Roux and Cline's limit $\Omega_{b'} \lesssim 0.8\,\Omega_{b}$ \cite{RouxCline2020}, only the helium interval has room: there T14b's Galaxy models pass at the Sun only with helium-star feedback \cite[Obs.~9.1]{ButlerT14b}, and Roux and Cline's limit, which covers the whole helium interval, is exceeded by a factor $1.24$. Nothing in the construction lowers the mirror's share below $0.995$ or changes its microphysics without a further exception to P2: the copy is exact and the tension T14 records stands. (The cold mirror itself is an exception to P2 at its version-1 reach; P2 as scoped leaves $\xi$ free, Theorem \emph{The tower charge is a sector label; the temperature is not}, and the cold mirror then needs a non-equilibrium descent, Proposition \emph{A cold mirror needs a non-equilibrium descent}.) Observation level, on imported galaxy modelling.
\end{conjecture}

\phantomsection\label{c:note.dark-tower-matter.setting-content-matter}
This section uses imported physics (gauge theory in three and five dimensions, \ref{c:rule2.imported-parity-anomaly} and \ref{c:rule2.imported-five-dimensional-chern-simons-shift}, and the published dark-plasma bounds) and decimal estimates, labelled; every number is reproduced by Part K of the script or by the named supplementary check. Its decimal estimates take T14's particle model (Fact \ref{c:dark-tower-matter.imported-unit-ball-rule}): the quanta's mass is the nucleon's and $d$ is this world's baryon spacing, $d_{b} \approx 0.98$ m. E alone fixes the mass density and leaves both the number of quanta and $d$ free, so a statement below that uses a number density or the value of $d$ is conditional on that model. The five-dimensional statements of Computation \ref{c:dark-tower-matter.five-dimensional-towers}(i), (ii) and (iv) need only $LT \gg 1$; the radiation caps of (iii) scale with $LT$.

\begin{theorem}[The phases of the three-dimensional towers]\label{c:dark-tower-matter.three-dimensional-phases}
\changedmark{BurntOrange}{changed since the deposit}
Imported: Polyakov confinement of a compact $U(1)$ with Maxwell dynamics \cite{Polyakov1977}; topological mass $e^{2}k/2\pi$ and the absence of confinement at nonzero Chern-Simons level \cite{DJT1982, AHPS1989}; flux attachment, with spin shifted by $q^{2}/2k$ \cite{Wilczek1982}. With the minimal bare level $|k| = \tfrac{1}{2}$ of Proposition \ref{c:dark-tower-matter.anomalies}(i), each factor is in one of two phases, set by whether the cell's mass sign is aligned with its bare level:
\begin{enumerate}
\item[(C)] effective level $0$: Polyakov-confining; the cells bind into neutral onia that annihilate at full strength into a lighter channel;
\item[(T)] effective level $\pm 1$: topologically massive, the photon gapped at $e^{2}/2\pi$; the unit-charge lines are transparent and the cells become effectively neutral and, on the standard signs, spin $0$ (interpretive).
\end{enumerate}
In phase (T) with $e^{2} \gg m$ the cells are stable by their gauge charges (at $r = 7$ also by the algebra's central unit, Proposition \ref{c:dark-tower-matter.residue-seven-unit}), and contact processes scale as $(m/e^{2})^{2}$. Assembly level. T14's Proposition 6.2 states that every compact $U(1)$ in three dimensions confines \cite[Prop.~6.2]{ButlerT14}; with the forced half-integral level that holds in phase (C) only. Both phases close T14's radiation channel: nothing light remains in phase (T), and even two massless $2+1$ dimensional photons at the photon temperature would give $\Delta\Neff$ of about $4 \times 10^{-3}$ on sheets (supplementary \file{CHECK\_COS}).
\end{theorem}
\begin{proof}
Follows from \ref{c:dark-tower-matter.anomalies}(i) (the minimal bare level), \ref{c:dark-tower-matter.three-dimensional-content} (the cells and their charges) and, for $r = 7$, \ref{c:dark-tower-matter.residue-seven-unit}, together with the imported results the source names (Polyakov confinement; the topological mass and the absence of confinement at nonzero level, after Deser, Jackiw and Templeton and Affleck, Harvey, Palla and Semenoff; flux attachment, after Wilczek), as the source states. Those imported results are not in the store.

\end{proof}

\begin{theorem}[Survival of sheet relics against annihilation]\label{c:dark-tower-matter.survival-sheets}
\changedmark{BurntOrange}{changed since the deposit}
Decimal. On a sheet the comoving number per unit area and the expansion rate both fall as $a^{-2}$ in radiation domination, so two-dimensional s-wave annihilation with $\sigma_{1}v = 2P/m$ ($P$ the absorption probability per encounter) depletes a symmetric population logarithmically, $1/N_{f} = 1/N_{0} + (2P/m)\ln(a_{f}/a_{i})/(a^{2}H)$. At the particle model's density a symmetric sheet relic keeps $99$ percent of its number only for $P$ below $1$ to $2 \times 10^{-16}$, the range spanning the number of relativistic species; at $P = 1$ a fraction of order $10^{-14}$ survives ($1.2 \times 10^{-14}$ with the inputs here, K3; T14 gives about $3 \times 10^{-14}$ for the same estimate \cite[Sec.~7.7]{ButlerT14}). In phase (T) this needs $e^{2}/m$ above $7 \times 10^{7}$ to $9 \times 10^{7}$, which also makes the relic collisionless in clusters (threshold $5.9 \times 10^{6}$, from the in-sheet phase-shift bound $1.7 \times 10^{-7}$ of T14b's cluster computation \cite[Comp.~3.3, Prem.~2.12]{ButlerT14b}) (K4). Phase (C) fails the particle model's density by itself.
\end{theorem}
\begin{proof}
Follows from the density of \ref{c:dark-tower-matter.imported-unit-ball-rule}, with the particle model's count, R1 (\ref{c:dark-matter.r1-equal-mass-per-quantum}) and R2 (\ref{c:dark-matter.r2-one-quantum-per-unit-ball}), for the number of quanta, the phases (T) and (C) of \ref{c:dark-tower-matter.three-dimensional-phases} and the in-sheet phase-shift bound of \ref{c:simulations.cluster-bound-two-dimensions}, as the source states; the decimal estimates are those the source cites to its script (K3, K4).

\end{proof}

\begin{theorem}[The five-dimensional towers in cosmology and halos]\label{c:dark-tower-matter.five-dimensional-towers}
\changedmark{BurntOrange}{changed since the deposit}
Decimal, with imports as stated. With the fifth direction one $d$ deep and co-expanding (R3), and $d = d_{b}$ (the particle model, R1 and R2; \cite[Prem.~7.4, Rem.~7.9, Rem.~7.14]{ButlerT14}). Under E alone the depth has no stated value, and the numbers below are conditional on the particle model's:
\begin{enumerate}
\item[(i)] $LT$ runs from about $390$ (at $g_{s} = 106.75$) to $1170$ today: the Kaluza-Klein spacing is always below $T/120$, so the towers are five-dimensional at the visible temperature and at their matter's mass scale (K5); at T14's small dark temperatures the dark bath itself is only a few Kaluza-Klein levels deep.
\item[(ii)] Five-dimensional gauge theory is infrared free; reducing to four dimensions gives a confinement scale below $\pi/L$ by a factor smaller than $e^{-250}$ even at strong coupling (K7), and the $1/r^{2}$ potential binds nothing unless the theory is within a factor of about eight of strong coupling at the quark mass (supplementary \file{D5C}, H3; \cite{CLP2000}). There are no dark hadrons and no dark atoms.
\item[(iii)] With the Kaluza-Klein tower populated, $\Delta\Neff < 0.3$ needs $T_{D}/T_{\gamma} < 0.085$ for the nine vectors of $\uu(3)$ ($0.13$ for one photon), not the $0.46$ to $0.51$ of T14's zero-mode formula \cite[Prop.~6.3]{ButlerT14}; a sector decoupled at a GeV would give about $670$ (K6). T14's own temperature estimates meet the corrected bound by orders.
\item[(iv)] Two-body scattering is collisionless by geometry, $\sigma_{\mathrm{eff}}/m = \sigma_{5}/(Lm) < 10^{-8}\ \mathrm{cm^{2}\,g^{-1}}$ even for charged and coloured species (supplementary \file{D5C}, H9).
\item[(v)] With no boundary term in the series the gauge field's natural condition at the ends of the fifth direction keeps a four-dimensional zero mode, so beyond about a metre the charged species form a neutral dark plasma. Published plasma-instability limits \cite{SSS2024, PIC2024}, rescaled to a GeV, lie six to ten orders below the coupling a freeze-out relic needs and about two to five orders below a deposit at a natural five-dimensional coupling (supplementary \file{D5C}, checks I3 and I4, and \file{CHECK\_COS}). Whether such instabilities transfer bulk momentum in a merger is not established \cite{Lasenby2020}: an exposure, not an exclusion.
\item[(vi)] If the neutral species $N$ (Proposition \ref{c:dark-tower-matter.five-dimensional-content}) carries the particle model's count as a deposit and the charged species annihilate at weak coupling, the tower carries no long-range charge and is collisionless whatever the zero modes do. Under Premise \ref{c:dark-tower-matter.u-transposition-rule} this branch is open to both five-dimensional towers.
\end{enumerate}
The relic is a symmetric gas of free spin-$\tfrac{1}{2}$ dark quarks and leptons, or of $N$ alone. The lightest species is stable by fermion parity, the lightest quark and lepton by species numbers broken only at the rung scale.
\end{theorem}
\begin{proof}
Follows from \ref{c:dark-tower-matter.gauge-algebras} (the nine vectors of $\uu(3)$) and \ref{c:dark-tower-matter.five-dimensional-content} (the neutral species $N$), with the fifth direction fixed by premise R3 (\ref{c:dark-matter.r3-one-cell-per-unshared-direction}), whose unit is this world's baryon spacing under R1 (\ref{c:dark-matter.r1-equal-mass-per-quantum}) and R2 (\ref{c:dark-matter.r2-one-quantum-per-unit-ball}), the particle model the source works in, and clause (vi) under Premise \ref{c:dark-tower-matter.u-transposition-rule}, as the source states. The decimal estimates are those the source cites to its script (K5 to K7) and to the supplementary script \file{D5C}; the imported results it names (five-dimensional strong coupling, plasma-instability limits) are not in the store.

\end{proof}

\begin{theorem}[Neither E nor the particle model has a mechanism in the odd towers]\label{c:dark-tower-matter.r1-has-no-mechanism-odd-towers}
\changedmark{BurntOrange}{changed since the deposit}
In this tower and the mirror a quantum's mass is a confinement scale, $m = c\Lambda$ \cite[Prop.~4.7]{ButlerT14}. In the three-dimensional towers the gauge algebra is abelian and the coupling dimensionful, so there is no dimensional transmutation and the cells' masses are bare parameters; the viable phase needs $m/e^{2}$ below $1.1$ to $1.5 \times 10^{-8}$, which nothing in the content protects (no chirality; parity broken by the forced Chern-Simons level). In the five-dimensional towers nothing confines (Computation \ref{c:dark-tower-matter.five-dimensional-towers}(ii)) and the quanta are elementary with bare masses. The odd towers carry no asymmetry, so nothing fixes their number either (Observation \ref{c:dark-tower-matter.population}). So nothing in the content fixes the masses or the densities that E equates across towers, for about $81$ percent of the dark matter; still less does it supply the particle model's equal masses. T14 already concedes that nothing sets an odd tower's quanta's mass to ours \cite[Rem.~7.8, Sec.~8]{ButlerT14}. What this note adds is the reason, and one correction: T14's Section 8 reads the odd towers' quanta as composites of sectors whose confinement scales the series does not derive, whereas here they are elementary in five dimensions and unconfined in the viable three-dimensional phase. Assembly level.
\end{theorem}
\begin{proof}
Follows from \ref{c:dark-matter.mirror-share-one-how-fragile} ($m = c\Lambda$ in this tower and the mirror), \ref{c:dark-tower-matter.gauge-algebras} (abelian in three dimensions), \ref{c:dark-tower-matter.three-dimensional-content} (no chirality), \ref{c:dark-tower-matter.anomalies} (the forced Chern-Simons level), \ref{c:dark-tower-matter.survival-sheets} (the viable phase), \ref{c:dark-tower-matter.five-dimensional-towers}(ii), the shares of \ref{c:dark-tower-matter.imported-unit-ball-rule} under Premise E (\ref{c:dark-matter.e-equal-mass-per-unit-ball}) and the population question of \ref{c:dark-tower-matter.population}, as the source states.

\end{proof}

\begin{theorem}[How the odd towers could be populated]\label{c:dark-tower-matter.population}
\changedmark{BurntOrange}{changed since the deposit}
The odd towers carry no asymmetry (Fact \ref{c:dark-tower-matter.imported-towers-their-types}). Their matter is therefore symmetric, and the number of quanta that fills E's density must be either a count deposited by the descent and protected from annihilation by kinematics or weakness (Computations \ref{c:dark-tower-matter.survival-sheets}, \ref{c:dark-tower-matter.five-dimensional-towers}(vi)), or, in the five-dimensional towers, a weak-coupling freeze-out with the coupling tied to the deposit temperature. Neither is derived. Assembly level.
\end{theorem}
\begin{proof}
Follows from \ref{c:dark-tower-matter.imported-towers-their-types} (no asymmetry in the odd towers), \ref{c:dark-tower-matter.survival-sheets} and \ref{c:dark-tower-matter.five-dimensional-towers}(vi), with the density of Premise E (\ref{c:dark-matter.e-equal-mass-per-unit-ball}) through \ref{c:dark-tower-matter.imported-unit-ball-rule}, as the source states; R2 (\ref{c:dark-matter.r2-one-quantum-per-unit-ball}) gives the count in the particle model.

\end{proof}

\section{From content to matter}

\begin{proposition}[The kinematics of the sterile towers]\label{c:dark-tower-matter.sterile-kinematics}
\changedmark{BurntOrange}{changed since the deposit}
\begin{enumerate}
\item[(i)] In $\R^{2,2}$ the orthogonal complement of every non-null vector is nondegenerate and indefinite, so every vector, of either sign of $Q$, is a sum of two null vectors $v = (\tfrac{1}{2}v + w) + (\tfrac{1}{2}v - w)$ with $w \perp v$, $Q(w) = -Q(v)/4$; and the null cone minus the origin is connected, so it has no invariant future half.
\item[(ii)] In $\R^{4,0}$ (and $\R^{0,4}$) the only null vector is zero: there is no massless constituent, no time and no evolution.
\end{enumerate}
\end{proposition}
\begin{proof}
\textit{(i) Sums of null vectors.}\enspace For non-null $v$, $v^{\perp}$ is three-dimensional, of signature $(1,2)$ if $Q(v) > 0$ and $(2,1)$ if $Q(v) < 0$, since the counts of \ref{c:synthesis.completion-recovers-the-real-clifford-algebras}(iii) add over $V = \R v \perp v^{\perp}$. Either way $Q$ is indefinite on it, so it takes the value $-Q(v)/4$ at some $w \perp v$, and then
\[
Q\bigl(\tfrac{1}{2}v \pm w\bigr) \;=\; \tfrac{1}{4}Q(v) + Q(w) \;=\; \tfrac{1}{4}Q(v) - \tfrac{1}{4}Q(v) \;=\; 0
\]
(J1, which also checks nondegeneracy and indefiniteness on random vectors, and J2, which exhibits sums of both signs).



\emph{(i) The cone is connected.}  In coordinates $(t_{1}, t_{2}; x_{1}, x_{2})$ with $Q = t_{1}^{2} + t_{2}^{2} - x_{1}^{2} - x_{2}^{2}$, the null cone minus the origin is
\[
\bigl\{\, |t| = |x| > 0 \,\bigr\} \;\cong\; S^{1} \times S^{1} \times \R_{>0},
\]
which is connected. In signature $(1,3)$ the same set has the two components $t_{0} > 0$ and $t_{0} < 0$. T16 proves the corresponding statement for the timelike sets \cite[Lem.~4.4]{ButlerT16}.



\emph{(ii)}  $Q$ is a sum of four squares, so $Q(v) = 0$ only at $v = 0$ (J3).

\end{proof}

\begin{remark}[The Kleinian algebra]\label{c:note.dark-tower-matter.kleinian-algebra}
At $r = 0$ the constituent field is $\R$ (Theorem \ref{c:dark-tower-matter.constituents-field}), and $\Cl^{0}(2,2) \cong \Mat_{2}(\R) \oplus \Mat_{2}(\R)$ has two inequivalent spinor ideals. There is no time direction, so no adjoint and no norm clause: a momentum bilinear in one element of each ideal is preserved by acting with $U$ on the elements of the first ideal and $U^{-T}$ on those of the second, for every $U \in GL(n,\R)$ (supplementary \file{STE}, checks C4, C5). Under H2 with this slot group the gauge algebra is the split form $\sll(2,\R) \oplus \sll(3,\R) \oplus \gl(1,\R)$, of dimension $12$, whose simple factors have indefinite Killing forms, signatures $(2,1)$ and $(5,3)$ (E5): there is no positive gauge-kinetic energy. T14's ``Kleinian copy of the Standard Model'' \cite[Cor.~3.11]{ButlerT14} needs a compact weak factor, which does not preserve that bilinear. Assembly level.
\end{remark}

\phantomsection\label{c:note.dark-tower-matter.setting-sterile-towers}
Both sterile towers are empty under T14's Premise P3, which is what the unit-ball rule's agreement with the observed ratio assumes: a populated $r = 0$ at share one would move the total to $6.36$, fifteen standard deviations from the observed $5.364 \pm 0.065$ \cite{Planck2018}.

\begin{conjecture}[The dark towers as built, against the data]\label{c:dark-tower-matter.towers-built}
\changedmark{BurntOrange}{changed since the deposit}
\begin{center}\small
\begin{tabular}{@{}p{3.6cm}p{10.2cm}@{}}
\toprule
test & status \\
\midrule
$\Delta\Neff < 0.3$ & met on every branch except a hot mirror ($\xi = 1$, $7.45$) and a five-dimensional sector decoupled at a GeV (about $670$) \\
cluster mergers & passes in T14b's runs only if the mirror is the only colliding component (Fact \ref{c:dark-tower-matter.imported-cluster-simulations}); that needs both three-dimensional towers in phase (T) at strong coupling with an unprotected light mass, and both five-dimensional towers collisionless (neutral-$N$ deposit, couplings weak near the unification scale, or ends the series does not supply). With zero modes and charged relics in both five-dimensional towers the colliding fraction would be $50.2$ percent, the fraction of T14b's failing mixed-count runs; with one, $34.4$ percent, above the $31.6$ percent run that already loses about half the subcluster's central mass. The extrapolation from fluids to a plasma is itself unproved \\
mirror bounds & marginal at best, in the helium interval with helium-star feedback, still $1.24$ above Roux and Cline; the helium interval itself needs an exception to P2 \\
R1 & a premise about bare masses in the odd towers; it holds only for the mirror \\
\bottomrule
\end{tabular}
\end{center}
No configuration of the five towers as built is compatible with every test at face value. The least exposed configuration needs four choices the series does not make (a P2 exception for the mirror's temperature, boosted mirror feedback, strong coupling with an unprotected light mass in the three-dimensional towers, and collisionless five-dimensional relics), and even with them it stays marginal, $1.24$ above Roux and Cline's limit. The P2 exception for the mirror's temperature is one at P2's version-1 reach; P2 as scoped leaves $\xi$ free (Theorem \emph{The tower charge is a sector label; the temperature is not}), and the cold mirror then needs a non-equilibrium descent (Proposition \emph{A cold mirror needs a non-equilibrium descent}). Observation level.
\end{conjecture}

\section{The sterile towers}

\phantomsection\label{c:note.dark-tower-matter.setting-what-note-depends-what-would}
\emph{Proved here, at proof level} (the two mass laws of Theorem \ref{c:dark-tower-matter.where-mass-begins} by exact symbolic verification)\emph{:} Theorem \ref{c:dark-tower-matter.constituents-field}, Lemma \ref{c:dark-tower-matter.adjoint-norm}, Theorems \ref{c:dark-tower-matter.where-mass-begins}, \ref{c:dark-tower-matter.ladder}, \ref{c:dark-tower-matter.third-pairs-unit-five-dimensions} and \ref{c:dark-tower-matter.mirror-algebra}, Proposition \ref{c:dark-tower-matter.pairs}, Proposition \ref{c:dark-tower-matter.residue-seven-unit}(i) to (iii) up to its last clause, Proposition \ref{c:dark-tower-matter.sterile-kinematics}, and the algebra in Theorem \ref{c:dark-tower-matter.gauge-algebras} and Propositions \ref{c:dark-tower-matter.three-dimensional-content}, \ref{c:dark-tower-matter.five-dimensional-content}, \ref{c:dark-tower-matter.anomalies}.

\emph{Assembly, under named premises:} that the algebras of Theorem \ref{c:dark-tower-matter.gauge-algebras} are gauge algebras (H2), that a constituent's momentum is $j$ of its spinor part off dimension four (H1), the constituent-level content (S), the transposition rule (U), including the compactness of the pair group (U2), which is assumed and not derived, and Proposition \ref{c:dark-tower-matter.mirrors-temperature-p2}, Observations \ref{c:dark-tower-matter.three-dimensional-phases}, \ref{c:dark-tower-matter.r1-has-no-mechanism-odd-towers}, \ref{c:dark-tower-matter.population}, Remark \ref{c:note.dark-tower-matter.kleinian-algebra}.

\emph{Decimal:} Computations \ref{c:dark-tower-matter.mirrors-composition}, \ref{c:dark-tower-matter.survival-sheets}, \ref{c:dark-tower-matter.five-dimensional-towers}, Observation \ref{c:dark-tower-matter.towers-built}.

\emph{Interpretive:} the mirror as parity image; the spin $0$ of the dressed three-dimensional cells.

\emph{What would refute the constructions.} A spin-irreducible constituent ideal whose spin commutant differs from Theorem \ref{c:dark-tower-matter.constituents-field}'s; a nonzero element of $S$ in dimension three or four with $\rev{\psi}\psi \neq 0$; an element of $\HH$ outside $\mathrm{span}\{1, I_{7}\}$ at $r = 3$ commuting with the $I_{7}$-built pair algebra. Each is a finite computation the script performs.

\emph{What would refute the physics.} A thin dark disk or bulge component excluding a mirror of $0.995\,\Omega_{b}$ in the helium interval with helium-star feedback refutes T14's mirror; a demonstration that dark plasma instabilities transfer bulk momentum in cluster mergers, with gauge zero modes in the five-dimensional towers, refutes the composition unless their relic is the neutral species; a measured $\Delta\Neff$ excess attributable to a dark sector would place a five-dimensional tower above $0.085$ of the photon temperature, outside T14's estimates.

\emph{Successor problems.} (1) Whether the cold mirror is a stated exception to P2, and where in T14's interval $\xi$ lies. (2) A footing for E's common length and equal mass per ball in the odd towers, where neither E nor the particle model has a mechanism; and, for T14, its Section 8's reading of the odd towers' quanta as confined composites, which Observation \ref{c:dark-tower-matter.r1-has-no-mechanism-odd-towers} contradicts. (3) The boundary data of the five-dimensional towers' fifth direction (zero modes or not) and their population mechanism. (4) The three-dimensional phase and coupling. (5) Hypothesis (S) off dimension four, the same missing theorem as in this tower. (6) The family count in the odd towers, where the maximal-composite bound of T13 \cite[Conj.~3.1]{ButlerT13} does not transpose. (7) The $r = 5$ unit: datum or field. (8) Whether the five-dimensional pair unit is internal (U3) or external, which decides between $\uu(3)$ and the Standard-Model algebra in both towers. (9) The $SU(3)$ and mixed Chern-Simons levels of the five-dimensional content.

\section{The dark sector as built, against the data}

\phantomsection\label{c:note.dark-tower-matter.setting-position-series}
This paper is T14c of the series, a companion to the dark-matter paper T14 \cite{ButlerT14}, downstream of its Theorem 3.4 (types by residue), Proposition 4.1 (superselection), Theorem 4.3 (the mirror copy), Theorem 3.10 (no orientation bit in odd dimension) and its unit-ball rule on Premises E and R3 (Section 7 of T14 Version 2.1). It runs the construction of the Standard-Model paper T12 \cite{ButlerT12} in each tower, using the mass formula of the mass paper T11 \cite{ButlerT11}, the rungs and canonical relations of the tower paper T3 \cite{ButlerT3}, and, for comparison only, the families of T13 \cite{ButlerT13} and the neutrino sector of T15 \cite{ButlerT15}. Its physical tests take the simulations of the companion note T14b \cite{ButlerT14b} as given; the companion note T14a \cite{ButlerT14a} is upstream of nothing here. The null-cone statement for signature $(2,2)$ is proved here (Proposition \ref{c:dark-tower-matter.sterile-kinematics}); the foundations paper T16 proves the corresponding statement for the timelike sets \cite{ButlerT16}. The capstone T17 \cite{ButlerT17} records T14's predictions; the section ``What this note depends on, and what would refute it'' lists what this note adds to them.

\phantomsection\label{c:note.dark-tower-matter.setting-machine-verification}
The script \file{verify\_T14c.py} \cite{ButlerVerifyT14c} imports the series library \file{GA.py} and checks the computational statements of Sections 3 to 10 in exact arithmetic: rational numbers, and a small polynomial ring over $\mathbb{Q}$ in which the identities of Lemma \ref{c:dark-tower-matter.adjoint-norm} and Theorems \ref{c:dark-tower-matter.where-mass-begins} and \ref{c:dark-tower-matter.ladder} are verified as polynomial identities in the spinors' coordinates rather than on samples. Part K alone is decimal and labelled. Without the supplementary scripts it runs $56$ checks, all passing:

\begin{center} \begin{tabular}{@{}lll@{}}
\toprule
Part & statements & checks \\
\midrule
A & Theorem \ref{c:dark-tower-matter.constituents-field} & A1 to A5 \\
B & Lemma \ref{c:dark-tower-matter.adjoint-norm}, Theorem \ref{c:dark-tower-matter.where-mass-begins} & B0 to B7 \\
C & Theorem \ref{c:dark-tower-matter.ladder} & C1, C2 \\
D & Proposition \ref{c:dark-tower-matter.pairs}, Theorem \ref{c:dark-tower-matter.third-pairs-unit-five-dimensions} & D1a, D1b, D2 to D12 \\
E & Theorem \ref{c:dark-tower-matter.gauge-algebras}, Remark \ref{c:note.dark-tower-matter.kleinian-algebra} & E1 to E5 \\
F & Propositions \ref{c:dark-tower-matter.three-dimensional-content}, \ref{c:dark-tower-matter.five-dimensional-content} & F1 to F4 \\
G & Proposition \ref{c:dark-tower-matter.anomalies} & G1, G2 \\
H & Proposition \ref{c:dark-tower-matter.residue-seven-unit} & H1 to H4 \\
I & Theorem \ref{c:dark-tower-matter.mirror-algebra} & I1 to I3 \\
J & Proposition \ref{c:dark-tower-matter.sterile-kinematics} & J1 to J3 \\
K & Proposition \ref{c:dark-tower-matter.mirrors-temperature-p2}, Computations \ref{c:dark-tower-matter.survival-sheets}, \ref{c:dark-tower-matter.five-dimensional-towers}; Fact \ref{c:dark-tower-matter.imported-unit-ball-rule} (decimal) & K1 to K7 \\
\bottomrule
\end{tabular}
\end{center}

Part L runs the supplementary scripts, which are deposited with the paper in its folder \file{T14c\_supplement}: the derivation scripts \file{MIR}, \file{D3R}, \file{D3C}, \file{D5C}, \file{D5Q}, \file{STE} (one per tower: $24$, $61$, $78$, $86$, $106$ and $68$ checks) and two independent checks, \file{CHECK\_ALG} ($50$ exact checks) and \file{CHECK\_COS} (decimal, ten checks on the numbers this paper cites), and requires each to be present and to run clean, twelve further checks ($68$ in all). The script needs \texttt{scipy} for Part K and \file{CHECK\_COS}, and takes about two and a half minutes. The supplementary scripts' docstrings refer to working notes of the run that produced them; those notes are not part of the deposit, and nothing here relies on them. The paper cites them for the decimal composition of the mirror (Computation \ref{c:dark-tower-matter.mirrors-composition}), the five-dimensional binding, scattering and plasma estimates (Computation \ref{c:dark-tower-matter.five-dimensional-towers}), the invariant forms of Corollary \ref{c:dark-tower-matter.same-kind-matter}, the branches of Remark \ref{c:note.dark-tower-matter.branches-break-u-what-they-give}, the Kleinian ladder of Remark \ref{c:note.dark-tower-matter.kleinian-algebra}, the transported weights of Proposition \ref{c:dark-tower-matter.three-dimensional-content}(iii) and the sheet radiation of Observation \ref{c:dark-tower-matter.three-dimensional-phases}. What no script checks: the statement after Theorem \ref{c:dark-tower-matter.gauge-algebras} about the mirror (T14's), the imported physics (Redlich's shift, Polyakov, topological mass, flux attachment, five-dimensional Chern-Simons quantisation, plasma instabilities) and the galaxy modelling of T14b.

\section{Dark matter and the dark towers}

\phantomsection\label{c:note.dark-tower-matter.setting-introduction}
The dark-matter paper T14 \cite{ButlerT14} derives eight parallel towers, one per residue $r = p - q \bmod 8$ of the real Clifford algebras, and finds that six can hold matter: this one ($r = 6$, signature $(1,3)$), its mirror ($r = 2$, $(3,1)$), two three-dimensional towers ($r = 1, 7$) and two five-dimensional towers ($r = 3, 5$). The towers touch only through the one shared metric. By its unit-ball rule the five dark towers hold, relative to our baryons, the shares $1$, $4/3$, $4/3$, $8/(3\pi)$ and $8/(3\pi)$, whose sum $(11\pi + 16)/(3\pi) = 5.364$ is the observed ratio of dark to visible matter. T14 proves the mirror a copy of this tower; for the four odd towers it says how much matter each holds and leaves open what that matter is.

This note answers that question by doing in each tower what the Standard-Model paper T12 \cite{ButlerT12} does in ours. T12's chain is: a constituent; the number of constituents at which mass begins; the mixings that preserve total momentum; the automorphisms of the smallest massive composite, identified with the gauge algebra; the content and its charges; then, with imported physics, hadrons, atoms and a stable relic. Every link is transposed here, and where the transposition meets an obstacle the resolution is taken from the way our tower resolves the same obstacle.

Everything is written in geometric algebra. A constituent is an element of a minimal left ideal of an even Clifford algebra; its momentum is the vector part of $\psi e_{t} \rev{\psi}$; its field of scalars is the corner algebra $f \Cl^{0} f$ with the adjoint $x^{\dagger} = e_{t}\rev{x}e_{t}^{-1}$. The mass laws are identities of the reversion pairing $\rev{\psi}_{1}\psi_{2}$.

\subsection*{Prior art}
\phantomsection\label{c:note.dark-tower-matter.setting-prior-art}
The real Clifford algebras, their periodicity and their spinor modules are classical (Cartan, Chevalley \cite{Chevalley1954}, Lounesto \cite{Lounesto2001}); the even-subalgebra isomorphisms $\Cl^{0}(p,q) \cong \Cl(p,q-1)$ and $\Cl^{0}(p,q) \cong \Cl^{0}(q,p)$ are standard \cite{Porteous1995, Lounesto2001}, and the division algebra of each even algebra can be read off the periodicity table. The spacetime-algebra current $j = \psi\gamma_{0}\rev{\psi}$ and its square $j^{2} = |\psi\rev{\psi}|^{2}$ are Hestenes's \cite{Hestenes1966, DoranLasenby2003}. Mirror matter and its cosmology have a long literature \cite{FootMirror2004, RouxCline2020}; the parity anomaly and topologically massive gauge fields in $2+1$ dimensions are Redlich's and Deser, Jackiw and Templeton's \cite{Redlich1984, DJT1982}; Polyakov confinement \cite{Polyakov1977}; five-dimensional gauge theory, its Chern-Simons terms and its strong-coupling scale are standard \cite{IMS1997, CLP2000}. The identification of a gauge algebra with the automorphisms of the minimal massive composite is T12's, and its global form and content are classical structure theory once that identification is made.

\subsection*{Contribution}
\phantomsection\label{c:note.dark-tower-matter.setting-contribution}
\begin{enumerate}
\item The constituent's field in tower $r$ is the division algebra of residue $r + 1$, at the conformal and the tower level alike (Theorem \ref{c:dark-tower-matter.constituents-field}, a corollary of the periodicity table): $\R$, $\C$, $\HH$ as the causal tower's dimension runs three, four, five. It is not the type T14 assigns by residue, which is the commutant of a different module (Remark \ref{c:note.dark-tower-matter.t14s-types-are-different-commutant}).
\item One identity of the reversion pairing gives where mass begins: at two constituents in three and four dimensions, where the pairing is alternating, and at one in five, where it is not (Theorem \ref{c:dark-tower-matter.where-mass-begins}). The momentum-preserving mixings are $O(n)$, $U(n)$, $Sp(n)$ (Theorem \ref{c:dark-tower-matter.ladder}).
\item The third ladder pair of the five-dimensional towers needs a complex unit. At $r = 3$ the algebra supplies one, and every other quaternionic unit fails to commute with the pair structure; at $r = 5$ the units form a sphere (Theorem \ref{c:dark-tower-matter.third-pairs-unit-five-dimensions}).
\item Under T12's identification premise and one transposition rule (Premise \ref{c:dark-tower-matter.u-transposition-rule}): $\so(2) \oplus \so(2)$ in both three-dimensional towers, $\uu(3)$ in both five-dimensional towers ($\su(2) \oplus \su(3) \oplus \uu(1)$ in both if the pair unit is external), the Standard Model in the mirror (Theorem \ref{c:dark-tower-matter.gauge-algebras} and, for the mirror, T14's copy theorem), with the content of Propositions \ref{c:dark-tower-matter.three-dimensional-content} and \ref{c:dark-tower-matter.five-dimensional-content}.
\item The towers of each odd dimension hold the same kind of matter (Corollary \ref{c:dark-tower-matter.same-kind-matter}); the classical isomorphism of the mirror's even algebras with this tower's is made explicit as a real map, which reverses the conformal pseudoscalar (Theorem \ref{c:dark-tower-matter.mirror-algebra}).
\item The physics: the matter of each tower in halos (the section ``From content to matter''); neither T14's premise E nor its particle model has a mechanism in the odd towers, whose quanta are elementary or unconfined with bare masses, not the confined composites T14's Section 8 describes (Observation \ref{c:dark-tower-matter.r1-has-no-mechanism-odd-towers}); a cold mirror is an exception to P2 (Proposition \ref{c:dark-tower-matter.mirrors-temperature-p2}); and the towers as built against the data (Observation \ref{c:dark-tower-matter.towers-built}).
\end{enumerate}

the table accompanying \ref{c:dark-tower-matter.imported-towers-their-types} and a figure of the original collect the result.

\centering  
\begin{tabular}{@{}llllll@{}}
\toprule
$r$ & tower & field & gauge algebra & elementary matter & in halos \\
\midrule
$6$ & $(1,3)$, ours & $\C$ & $\su(3) \oplus \su(2) \oplus \uu(1)$ & quarks, leptons & atoms, stars \\
$2$ & $(3,1)$, mirror & $\C$ & the same & a copy & dissipative gas \\
$1$ & $(2,1)$, sheets & $\R$ & $\so(2) \oplus \so(2)$ & cells $A$, $B$ & collisionless in one phase \\
$7$ & $(1,2)$, sheets & $\R$ & $\so(2) \oplus \so(2)$ & cells $A$, $B$ & collisionless in one phase \\
$3$ & $(4,1)$ & $\HH$ & $\uu(3)$\,$^{*}$ & $U, D, E, N$, vector-like & free fermions \\
$5$ & $(1,4)$ & $\HH$ & $\uu(3)$\,$^{*}$ & $U, D, E, N$, vector-like & free fermions \\
$0$ & $(2,2)$ & $\R$ & $\sll(2,\R) \oplus \sll(3,\R) \oplus \gl(1,\R)$ & none (no stable vacuum) & none \\
$4$ & $(4,0)$, $(0,4)$ & $\HH$ & (no time) & none & none \\
\bottomrule
\end{tabular}

\emph{The eight towers' matter. Fields: Theorem \ref{c:dark-tower-matter.constituents-field}. Gauge algebras: Theorem \ref{c:dark-tower-matter.gauge-algebras}, under H2 and Premise \ref{c:dark-tower-matter.u-transposition-rule}; $^{*}$with an external pair unit, $\su(2) \oplus \su(3) \oplus \uu(1)$ in both (Remark \ref{c:note.dark-tower-matter.branches-break-u-what-they-give}); $r = 0$: Remark \ref{c:note.dark-tower-matter.kleinian-algebra}. Content: Propositions \ref{c:dark-tower-matter.three-dimensional-content}, \ref{c:dark-tower-matter.five-dimensional-content}. In halos: the section ``From content to matter''.}

\section{Results used from elsewhere}
The proofs above use the following results of the series, stated briefly; each is proved in the paper named.

\begin{fact}[The flag bit exists exactly at tower dimension $0 \bmod 4$]\label{c:dark-matter.flag-bit-exists-exactly-tower-dimension}
The choice between the two halves of a Witt pair is a bit of the connected conformal spin group exactly when the tower dimension is $0 \bmod 4$, so, under Premise P3 (nearest-dimension representatives), the even residues can be oriented, and the odd residues never can.
\emph{Source:} T14, Thm 3.10 (\href{https://doi.org/10.5281/zenodo.22722953}{doi:10.5281/zenodo.22722953}).
\end{fact}

\begin{fact}[The mirror share is one, and how fragile that is]\label{c:dark-matter.mirror-share-one-how-fragile}
Under the hypothesis of equal conversion (the mirror's ratio of baryon number to $B - L$ equals this tower's, which is not shown), the mirror share is exactly one because the mirror copies this tower (which discharges the democracy premise for it), but at one loop, over the ceiling range of Hypothesis (U), one extra Dirac flavour, or a ceiling coupling difference above about $0.1$, would move it far from one, and nothing follows for the odd towers.
\emph{Source:} T14, Prop 4.7 (\href{https://doi.org/10.5281/zenodo.22722953}{doi:10.5281/zenodo.22722953}).
\end{fact}

\begin{fact}[The mirror tower is a copy of this one]\label{c:dark-matter.mirror-tower-copy-one}
Every input of the Standard-Model chain is invariant under signature reversal, so the mirror tower carries a copy of the derived Standard Model and, with a common ceiling $M$ and equal couplings there (the common-ceiling-coupling premise) and under the single frozen point and indifference (which give the equal asymmetry per mode on the charge reading, an equality of $B - L$) and under the hypothesis of equal conversion (the mirror's ratio of baryon number to $B - L$ equals this tower's, which is not shown), $\Omega_{\mathrm{mirror}} = \Omega_{b}$ to half a percent; indifference at its first reach would also make the mirror as hot as ours, and as scoped it leaves the temperature free.
\emph{Source:} T14, Thm 4.3 (\href{https://doi.org/10.5281/zenodo.22722953}{doi:10.5281/zenodo.22722953}).
\end{fact}

\begin{fact}[Seeding by gravity: the floor on a dark tower's temperature]\label{c:dark-matter.seeding-gravity-floor-dark-towers-temperature}
Gravitational pair production from the visible plasma, at the rate derived in the store (coefficient $c$, counting events of two quanta, each carrying a mean energy of about $7.9\,T$ computed from the rate integral), puts a floor $\xi_{\mathrm{floor}} \propto c^{1/4}T_{\mathrm{desc}}^{3/4}$ on the temperature ratio of a dark tower whose light species thermalise among themselves, $1.23 \times 10^{-3}\,c^{1/4}(16/g_{d})^{1/4}$ at the ceiling reheating.
\emph{Source:} T14, Prop A.2 (\href{https://doi.org/10.5281/zenodo.22722953}{doi:10.5281/zenodo.22722953}).
\end{fact}

\begin{fact}[The stored asymmetry is a charge, not a number]\label{c:dark-matter.stored-asymmetry-charge-not-number}
Under the inherited-asymmetry classification, with the number $N$ (which fixes a count of registered units) common to the towers under P1, an oriented tower's asymmetry is the temperature-independent charge $Q = \epsilon N$, named as $B - L$ (not the residue label, not baryon number), so the mirror carries the same $B - L$ at any mirror temperature, and the mirror share $f_{2} = 1$ holds if the ratio of baryon number to $B - L$ is the same in the two towers.
\emph{Source:} T14, Prop 4.6 (\href{https://doi.org/10.5281/zenodo.22722953}{doi:10.5281/zenodo.22722953}).
\end{fact}

\begin{fact}[Superselected, hence dark and exactly stable]\label{c:dark-matter.superselected-hence-dark-exactly-stable}
Under P1, the sector decomposition and the wall-sequence theorem's hypotheses (compact base, good couplings), the tower charge is exactly superselected, so dark-tower matter is exactly stable; that it only gravitates, at first order, is the series' acoustic no-go's, within the s-wave current algebra and conditional on the Indifference premise and premise (Q).
\emph{Source:} T14, Prop 4.1 (\href{https://doi.org/10.5281/zenodo.22722953}{doi:10.5281/zenodo.22722953}).
\end{fact}

\begin{fact}[Towers are residues, and the single-clock filter]\label{c:dark-matter.towers-are-residues-single-clock-filter}
The residue $p - q \pmod 8$ is invariant along each tower, and under the nearest-dimension premise P3, with P2 reading the time-sign label as indexical, exactly the residues $0$ and $4$ lack a single time direction, leaving six towers that admit matter (the two-time exclusion imports the standard instability).
\emph{Source:} T14, Prop 3.1 (\href{https://doi.org/10.5281/zenodo.22722953}{doi:10.5281/zenodo.22722953}).
\end{fact}

\begin{fact}[The two families are the two chiralities]\label{c:dark-matter.two-families-are-two-chiralities}
In even dimension $2m$ maximal isotropic subspaces fall into exactly two families, labelled by the chirality of their pure spinors, with a Witt pair in opposite families exactly when $m$ is odd; in odd dimension there is one family.
\emph{Source:} T14, Prop 3.9 (\href{https://doi.org/10.5281/zenodo.22722953}{doi:10.5281/zenodo.22722953}).
\end{fact}

\begin{fact}[Type by residue, and what each tower can carry intrinsically]\label{c:dark-matter.type-residue-what-each-tower-can}
A tower's spinor type (real, complex or quaternionic, with one or two modules) is a function of its residue alone, and it fixes the intrinsic symmetry group, $\Z_{2}$, $U(1)$ or $Sp(1)$, that each tower's ground algebra carries.
\emph{Source:} T14, Thm 3.4 (\href{https://doi.org/10.5281/zenodo.22722953}{doi:10.5281/zenodo.22722953}).
\end{fact}

\begin{fact}[The unit-ball rule]\label{c:dark-matter.unit-ball-rule}
Under Premises E (equal mass per unit ball) and R3, with the odd towers populated and residues 0 and 4 empty (both registered premises), a tower with $D$-dimensional space has share $V_{3}/V_{D}$, and the dark-to-baryon ratio is $(11\pi + 16)/(3\pi) \approx 5.3643$, against the observed $5.364 \pm 0.065$; the quanta's masses and numbers do not enter, E with R3 is equivalent to the shares tower by tower with equal thicknesses, and the agreement is built in by fitting, not predicted.
\emph{Source:} T14, Prop 7.6 (\href{https://doi.org/10.5281/zenodo.22722953}{doi:10.5281/zenodo.22722953}).
\end{fact}

\begin{fact}[Mass as symplectic area: the two-spinor formula]\label{c:mass.mass-symplectic-area-two-twistor-formula}
Two conformal spinors with momentum $\pi^{1}\bar\pi^{1} + \pi^{2}\bar\pi^{2}$ have $m^{2} = 2|\br{\pi^{1}\pi^{2}}|^{2}$, the squared symplectic area of their spinor parts.
\emph{Source:} T11, Prop 4.1 (\href{https://doi.org/10.5281/zenodo.22722565}{doi:10.5281/zenodo.22722565}).
\end{fact}

\begin{fact}[The sign of $\omega^2$]\label{c:primer.sign}
The pseudoscalar of $\Cl(p,q)$ satisfies $\omega^{2} = (-1)^{(p-q)(p-q-1)/2}$, so $\omega^{2} = +1$ exactly when $p - q \equiv 0$ or $1 \pmod 4$.
\emph{Source:} T0, Lem 9.5 (\href{https://doi.org/10.5281/zenodo.22714501}{doi:10.5281/zenodo.22714501}).
\end{fact}

\begin{fact}[Imported: the five-dimensional Chern-Simons shift]\label{c:rule2.imported-five-dimensional-chern-simons-shift}
Integrating out a massive five-dimensional Dirac fermion shifts the Chern-Simons coefficient by $\tfrac12\sgn(m)$, a parity-odd term present for one species and absent for a cancelling pair.
\emph{Source:} E. Witten, Phase transitions in M-theory and F-theory, Nucl. Phys. B 471 (1996) 195--216.; K. Intriligator, D. R. Morrison, N. Seiberg, Five-dimensional supersymmetric gauge theories and degenerations of Calabi-Yau spaces, Nucl. Phys. B 497 (1997) 56--100..
\end{fact}

\begin{fact}[Imported: the parity anomaly]\label{c:rule2.imported-parity-anomaly}
A single massive two-component Dirac fermion in $2+1$ dimensions induces a parity-odd Chern-Simons term at level $\tfrac12\sgn(m)$, which cancels for an equal-mass pair in the two inequivalent representations.
\emph{Source:} A. J. Niemi, G. W. Semenoff, Axial-anomaly-induced fermion fractionization and effective gauge-theory actions in odd-dimensional space-times, Phys. Rev. Lett. 51 (1983) 2077--2080.; A. N. Redlich, Gauge noninvariance and parity nonconservation of three-dimensional fermions, Phys. Rev. Lett. 52 (1984) 18--21; Parity violation and gauge noninvariance of the effective gauge field action in three dimensions, Phys. Rev. D 29 (1984) 2366--2374.; L. Alvarez-Gaum'e, E. Witten, Gravitational anomalies, Nucl. Phys. B 234 (1984) 269--330..
\end{fact}

\begin{fact}[The cluster bound in two dimensions]\label{c:simulations.cluster-bound-two-dimensions}
For nucleon-mass sheet quanta at cluster speeds, assuming premises R1 to R3 for the core spacing, the cluster bound of $1\ \cmg$ needs an in-sheet cross-section below $4.8 \times 10^{-12}$ fm, which requires either an impossibly small size or an s-wave phase shift below $1.7 \times 10^{-7}$.
\emph{Source:} T14b, Comp 3.3 (\href{https://doi.org/10.5281/zenodo.23234576}{doi:10.5281/zenodo.23234576}).
\end{fact}

\begin{fact}[The gauge algebra of the minimal massive composite]\label{c:standard-model.gauge-algebra-minimal-massive-composite}
The automorphism algebra of the minimal massive hierarchy $\C^{2} \otimes \C^{3}$ is $\su(2) \oplus \su(3) \oplus \uu(1)$, exactly the Standard Model's, while that of $\C^{3} \otimes \C^{3}$ is seventeen-dimensional.
\emph{Source:} T12, Thm 4.1 (\href{https://doi.org/10.5281/zenodo.22722669}{doi:10.5281/zenodo.22722669}).
\end{fact}

\begin{fact}[The global form: the quotient determined]\label{c:standard-model.global-form-quotient-determined}
The global gauge group of the minimal massive composite is $S(U(2) \times U(3)) \cong (SU(2) \times SU(3) \times U(1))/\mathbb{Z}_{6}$, the smallest of the Standard Model's global forms.
\emph{Source:} T12, Thm 4.2 (\href{https://doi.org/10.5281/zenodo.22722669}{doi:10.5281/zenodo.22722669}).
\end{fact}

\begin{fact}[Hypercharge from the global form, pair counting, and consistency]\label{c:standard-model.hypercharge-global-form-pair-counting-consistency}
Under Hypothesis (S) for the doublets, the composite's states carry the observed hypercharges of one generation, which satisfy all four anomaly conditions and are the unique solution for that content with $Y_{Q} = \tfrac16$.
\emph{Source:} T12, Thm 5.5 (\href{https://doi.org/10.5281/zenodo.22722669}{doi:10.5281/zenodo.22722669}).
\end{fact}

\begin{fact}[Imported: the pair systems and the two chirality halves]\label{c:standard-model.imported-pair-systems-two-chirality-halves}
Imported: the Fock module splits into the four-dimensional halves $\CS$ and $\CSm$, and every nonzero conformal spinor is a pure spinor, the Fock vacuum of its own three ladder pairs.
\emph{Source:} T12, Fact 2.2 (\href{https://doi.org/10.5281/zenodo.22722669}{doi:10.5281/zenodo.22722669}).
\end{fact}

\begin{fact}[The internal symmetry ladder]\label{c:standard-model.internal-symmetry-ladder}
The mixings of a system of $n$ conformal spinors that preserve total momentum and total helicity form exactly $U(n)$, giving the ladder $U(1)$, $U(2)$, $U(3)$.
\emph{Source:} T12, Prop 3.1 (\href{https://doi.org/10.5281/zenodo.22722669}{doi:10.5281/zenodo.22722669}).
\end{fact}

\begin{fact}[The pair count in dimension $n$]\label{c:standard-model.pair-count-dimension}
For even spacetime dimension $n$ a maximal isotropic subspace of $V_{\C}$ has dimension $(n+2)/2$, so the pair system of a conformal spinor has $(n+2)/2$ pairs, three in dimension four.
\emph{Source:} T12, Lem 2.12 (\href{https://doi.org/10.5281/zenodo.22722669}{doi:10.5281/zenodo.22722669}).
\end{fact}

\begin{fact}[Quarks and leptons: the two-level fermion content]\label{c:standard-model.quarks-leptons-two-level-fermion-content}
The pair-level cells carry the quark doublet $(2,3)_{1/6}$ unconditionally, and under Hypothesis (S) one slot with the pairs full gives the lepton doublet $(2,1)_{-1/2}$.
\emph{Source:} T12, Prop 5.3 (\href{https://doi.org/10.5281/zenodo.22722669}{doi:10.5281/zenodo.22722669}).
\end{fact}

\begin{fact}[The ten is the transported ambient space, and its doublets are the rungs]\label{c:standard-model.ten-transported-ambient-space-doublets-are}
The ten vectors of the transported ambient space decompose as $(2,1)_{1/2} \oplus (2,1)_{-1/2} \oplus (1,3)_{-1/3} \oplus (1,\bar{3})_{1/3}$, the rung directions carrying the Higgs doublet's representation.
\emph{Source:} T12, Prop 9.22 (\href{https://doi.org/10.5281/zenodo.22722669}{doi:10.5281/zenodo.22722669}).
\end{fact}

\begin{fact}[The three involutions on each grade]\label{c:std-algebra.involution-properties}
On $\Cl^{k}$: $\hat{x} = (-1)^{k}x$, $\tilde{x} = (-1)^{k(k-1)/2}x$, $\bar{x} = (-1)^{k(k+1)/2}x$; the first is an automorphism, the others anti-automorphisms.
\emph{Source:} the algebra core of the Physics Graph, proved there with an exact check; standard, see P. Lounesto, Clifford Algebras and Spinors, 2nd ed., London Mathematical Society Lecture Note Series 286, Cambridge University Press (2001), Ch. 14 (definitions of the Clifford algebra).
\end{fact}

\begin{fact}[How the pseudoscalar commutes]\label{c:std-algebra.pseudoscalar-commutation}
$I X_k = (-1)^{k(n-1)} X_k I$; $I\tilde{I} = \prod_i Q(e_i)$, and $I^{2} = 0$ in a degenerate algebra.
\emph{Source:} the algebra core of the Physics Graph, proved there with an exact check; standard, see C. Doran and A. Lasenby, Geometric Algebra for Physicists, Cambridge University Press (2003), Ch. 4 (foundations of geometric algebra).
\end{fact}

\begin{fact}[Completion recovers the real Clifford algebras; Sylvester's law of inertia in GA]\label{c:synthesis.completion-recovers-the-real-clifford-algebras}
If the coefficients form a complete ordered field (premise (F4)), every real quadratic space has an anticommuting basis of squares $1, -1, 0$ with invariant counts (Sylvester, derived in GA; Euclidean suffices) and every real Clifford algebra is $\Cl_{\Q}(p,q,r) \otimes \R$, inside $\Cl_{\Q}(p{+}r,q{+}r) \otimes \R$; over $\Q$ square classes remain ($\langle 2, 3\rangle$ is a division algebra).
\emph{Source:} T0a, Thm 6.1 (\href{https://doi.org/10.5281/zenodo.24402873}{doi:10.5281/zenodo.24402873}).
\end{fact}

\begin{fact}[The ladder pairs satisfy the CAR, as vectors]\label{c:tower.ladder-pairs-satisfy-car-vectors}
In the $k$-th tower algebra the $k$ adjoined null pairs satisfy the canonical anticommutation relations of $k$ fermionic modes exactly, as vectors, and each new rung adds one mode without altering the existing relations.
\emph{Source:} T3, Prop 4.1 (\href{https://doi.org/10.5281/zenodo.22721556}{doi:10.5281/zenodo.22721556}).
\end{fact}

\begin{fact}[Rigidity; the clock is frozen]\label{c:tower.rigidity-clock-frozen}
Every floor of the tower shares the ground floor's charge, division-ring type, centre, simplicity class, Morita class and obstruction taxonomy, so the physical tower $\Cl(1{+}k, 3{+}k) \cong \Mat_{2^{k+1}}(\HH)$ is frozen at clock position $6$.
\emph{Source:} T3, Thm 3.6 (\href{https://doi.org/10.5281/zenodo.22721556}{doi:10.5281/zenodo.22721556}).
\end{fact}

\begin{refutation}[The raw representation content matches the fermion table entry by entry]\label{c:refuted.the-raw-representation-content-matches-the-fermi}
Refuted: the raw representation content does not match the fermion table entry by entry; the disagreements are composites, the generators are the content, and the lepton lies one level higher under the slot-statistics hypothesis.
\end{refutation}

\begin{refutation}[No Standard Model algebra on the three-spinor hierarchy]\label{c:refuted.the-three-twistor-hierarchy-carries-the-standard}
Refuted: the hierarchy of three conformal spinors $\C^{3}\otimes\C^{3}$ has a seventeen-dimensional automorphism algebra, not the Standard Model's twelve, which the minimal massive composite $\C^{2}\otimes\C^{3}$ carries instead.
\end{refutation}

\begin{prediction}[Direct detection stays null]\label{c:pred.dark-matter-direct-detection-null}
Conditionally on the single-frozen-point premise, and on the series' acoustic no-go (s-wave current algebra, Indifference premise, premise (Q)), direct-detection searches should stay null permanently, because dark matter couples to ordinary matter through gravity alone at first order.
\emph{Kill condition:} Any confirmed dark-matter detection above the gravitational floor.
\end{prediction}

\begin{prediction}[No indirect dark-sector signals]\label{c:pred.no-indirect-dark-signals}
Conditionally on the single-frozen-point premise, and on the series' acoustic no-go (s-wave current algebra, Indifference premise, premise (Q)), indirect-detection surveys should see no dark-sector decay products, and no annihilation products above the gravitationally suppressed rate of a charge-asymmetric relic.
\emph{Kill condition:} Any confirmed indirect signal.
\end{prediction}

\begin{prediction}[No fractionally charged colour singlet]\label{c:pred.z6-charge-lattice}
The $\mathbb{Z}_{6}$ charge lattice of the derived gauge group forbids any fractionally charged colour singlet.
\emph{Kill condition:} One confirmed fractionally charged colour-singlet particle.
\end{prediction}

\section*{Changes since the deposited version}
\addcontentsline{toc}{section}{Changes since the deposited version}
Since the deposited version 1.0: 28 results changed, 2 reworded (wording, notation or citations only) and 38 cited results added. The comparison is with the source of version 1.0 on Zenodo, \href{https://doi.org/10.5281/zenodo.24403592}{doi:10.5281/zenodo.24403592}. A result counts as revised when the words or mathematics of its statement or proof differ from that source; a citation added in a parenthesis does not count. A result the paper cites from elsewhere is compared with the deposit of the paper that proves it. Each revision is classed as corrected, changed or reworded; the classes were assigned by a review of each difference and the store's record of why it was made, and a revision not yet reviewed is counted as changed.
\subsection*{Changed}
The content changed by a later decision, not because the deposited version was wrong.
\begin{itemize}\setlength{\itemsep}{2pt}
\item Postulate~\ref{c:dark-tower-matter.h1-momentum} (\emph{Premise H1: a constituent's momentum}), Premise 2.10 in the deposit: statement. \emph{Store notes.} \emph{2026-11-10:} Maintainer item M4 of 10 November 2026: a note at every level points to synthesis.constituent-momentum-is-the-unique-spin-covariant-quadratic-vector-on-the-spinor-ideal (the formula is forced up to a real factor given quadraticity and covariance under the connected spin group, in the causal towers of dimension three to five) and lists what stays premised [...]
\item Postulate~\ref{c:dark-tower-matter.u-transposition-rule} (\emph{Premise U: the transposition rule (a definition)}), Premise 2.13 in the deposit: statement. \emph{Store notes.} \emph{2026-11-10:} Leon's decision 10 of 10 November 2026: Premise U is typed as a definition. A clean re-typing as an object is not available: an object needs a new id (obj.*, this id then retired), cannot be the target of an assumes edge, and as a uses target raises W209, so the five results that assume U would lose the edge (one of them approved by Leon at acceptance) and [...] \emph{2026-11-20:} The formal level restated: T14c deposited as version 1.0 on 20 November 2026 (10.5281/zenodo.24403592), after an independent review: the norm clause of (U2) stated as part of the definition, with what rests on it. The technical level restated: as the formal level.
\item Assembled~\ref{c:dark-tower-matter.imported-towers-their-types} (\emph{Imported: the towers and their residues}), Fact 2.1 in the deposit: statement and proof.
\item Assembled~\ref{c:dark-tower-matter.imported-unit-ball-rule} (\emph{Imported: the unit-ball shares}), Fact 2.4 in the deposit: proof. \emph{Store notes.} \emph{2026-11-20:} The formal level restated: T14c deposited as version 1.0 on 20 November 2026 (10.5281/zenodo.24403592), after an independent review: T14's rule stated on its premises E and R3, the particle model R1 and R2 implying E and not used by the rule (T14 version 2.1); edge to E added. The technical level restated: as the formal level. The lay level restated: as the [...]
\item Fact~\ref{c:dark-tower-matter.imported-cluster-simulations} (\emph{Imported: the cluster merger simulations}), Fact 2.5 in the deposit: statement. \emph{Store notes.} \emph{2026-10-21:} Series import of the T14b merger runs, linked by lineage to the merger observation and from the towers-built observation; no proof uses it and no definition bears on it.
\item Theorem~\ref{c:dark-tower-matter.constituents-field} (\emph{The constituent's field in each tower}), Thm 3.2 in the deposit: statement and proof.
\item Lemma~\ref{c:dark-tower-matter.adjoint-norm} (\emph{The adjoint and the norm of the field}), Lem 3.4 in the deposit: proof.
\item Theorem~\ref{c:dark-tower-matter.where-mass-begins} (\emph{Where mass begins in each causal tower}), Thm 3.5 in the deposit: proof.
\item Theorem~\ref{c:dark-tower-matter.ladder} (\emph{The momentum-preserving mixings of constituents}), Thm 3.7 in the deposit: statement and proof.
\item Proposition~\ref{c:dark-tower-matter.pairs} (\emph{The ladder pairs of the causal towers}), Prop 4.1 in the deposit: proof.
\item Theorem~\ref{c:dark-tower-matter.third-pairs-unit-five-dimensions} (\emph{The third pair's unit in five dimensions}), Thm 4.2 in the deposit: statement and proof. \emph{Store notes.} \emph{2026-11-07:} Notation: I\_7 here is the pseudoscalar of Cl(5,2), the r = 3 tower's conformal algebra, while the conventions' I\_7 = I\_6 e is the unit of Cl(3,4); clause (i) now says which is meant. No change of mathematics.
\item Theorem~\ref{c:dark-tower-matter.gauge-algebras} (\emph{The gauge algebras of the odd towers}), Thm 5.1 in the deposit: proof.
\item Proposition~\ref{c:dark-tower-matter.three-dimensional-content} (\emph{The three-dimensional towers' content}), Prop 6.1 in the deposit: proof.
\item Remark~\ref{c:note.dark-tower-matter.neutral-row} (\emph{The neutral row}), Rem 6.2 in the deposit: statement. \emph{Store notes.} \emph{2026-11-07:} Unit principle: the note, written over span\{1, I\_4\} by ruling R3, now says that the five-dimensional towers' pairs are built with the hierarchy's unit u (I\_7 of Cl(5,2) at r = 3, a chosen quaternionic unit at r = 5) and that I\_4 is the multiplicity spaces' unit, not identified with u. No change of mathematics. \emph{2026-11-07:} By Leon's ruling R3: the central generator of u(1) + u(k) is (k I\_4, -I\_4 1\_k), the slot and pair multiplicity spaces C\^{}1 + C\^{}k read over span\{1, I\_4\} on the tower's floors (stated once; notational); no bare i remains, so the earlier declaration is removed.
\item Proposition~\ref{c:dark-tower-matter.five-dimensional-content} (\emph{The five-dimensional towers' content}), Prop 6.3 in the deposit: proof.
\item Proposition~\ref{c:dark-tower-matter.anomalies} (\emph{Chern-Simons levels of the odd towers' content}), Prop 6.5 in the deposit: statement and proof. \emph{Store notes.} \emph{2026-11-20:} T14c deposited as version 1.0 on 20 November 2026 (10.5281/zenodo.24403592), after an independent review: the legacy tags follow the revised verify\_T14c.py, whose checks G1 and G2 were merged into G1 (the old G3 is now G2).
\item Corollary~\ref{c:dark-tower-matter.same-kind-matter} (\emph{Paired towers hold the same kind of matter}), Cor 7.2 in the deposit: proof. \emph{Store notes.} \emph{2026-10-21:} The proof assembles rows from cited results, so U and H2 are given a closing role sentence.
\item Theorem~\ref{c:dark-tower-matter.mirror-algebra} (\emph{The mirror's even algebra as a real isomorphism}), Thm 8.1 in the deposit: proof.
\item Proposition~\ref{c:dark-tower-matter.mirrors-temperature-p2} (\emph{The mirror's temperature and premise P2}), Prop 8.2 in the deposit: statement and proof. \emph{Store notes.} \emph{2026-11-07:} Pin review: dark-matter.mirror-tower-copy-one repinned 8 -> 9 after reviewing its C1 restatement (clause (iv): the common ceiling and coupling are now Premise synthesis.common-ceiling-coupling-premise, P2 equating labels and not couplings; the share of one rests also on that premise; the caveat is reworded for P2's version-1 reach), clauses (i) to (iii) and [...] \emph{2026-11-07:} G1 (DECISIONS; R05 item 2): the cold-mirror exception now says that it is an exception to P2 at its version-1 reach, that P2 as scoped (version 2, labels not states) leaves xi free (synthesis.tower-charge-sector-label-temperature-not), and that the cold mirror then needs a non-equilibrium descent (synthesis.cold-mirror-needs-non-equilibrium-descent); P2 is [...]
\item Theorem~\ref{c:dark-tower-matter.mirrors-composition} (\emph{The mirror's composition across its temperature range}), Comp 8.3 in the deposit: statement and proof. \emph{Store notes.} \emph{2026-11-18:} , at Leon's decision of 18 November 2026, 14:04 (adopt the computed mean energy per creation event in the gravitational floor). Quoted numbers only: the present floor 5.2e-4 -> 6.6e-4 at every level and in the proof text; the floor value of eta' 4.4 -> 2.1. No number of clauses (i) to (iii) changed. Clause (i) remains stated at xi = 2.8 to 3.2e-4 (the floor [...] \emph{2026-11-18:} , at Leon's decision of 18 November 2026 (restate the numbers that used the first published production coefficients). The supplementary scripts behind clause (i) are not in the lab, so (i) could not be recomputed at the present floor 5.2e-4; it is restated as holding at xi = 2.8 to 3.2e-4, where it was computed, which is the floor for a reheating below the [...]
\item Conjecture~\ref{c:dark-tower-matter.mirror-against-record} (\emph{The mirror against the astrophysical record}), Obs 8.4 in the deposit: statement. \emph{Store notes.} \emph{2026-10-21:} Terminal observation on imported galaxy modelling; lineage now includes the imported Milky Way bounds (Widmark, Roux and Cline) it confronts. \emph{2026-11-07:} G1 (DECISIONS; R05 item 2): the cold-mirror exception now says that it is an exception to P2 at its version-1 reach, that P2 as scoped (version 2, labels not states) leaves xi free (synthesis.tower-charge-sector-label-temperature-not), and that the cold mirror then needs a non-equilibrium descent (synthesis.cold-mirror-needs-non-equilibrium-descent); P2 is [...]
\item Theorem~\ref{c:dark-tower-matter.three-dimensional-phases} (\emph{The phases of the three-dimensional towers}), Obs 9.1 in the deposit: statement and proof. \emph{Store notes.} \emph{2026-10-21:} Lineage supersedes added (decision 1, CONTRADICTIONS\_STRUCTURE H3): T14 Prop 6.2(2) said every compact U(1) in three dimensions confines by monopoles; with the half-integral level of Prop 6.5(i) that holds in phase (C) only. The audit's clause (2) and conclusion are narrowed and cite this result and the anomaly proposition. Lineage is unpinned in the store [...]
\item Theorem~\ref{c:dark-tower-matter.survival-sheets} (\emph{Survival of sheet relics against annihilation}), Comp 9.2 in the deposit: proof. \emph{Store notes.} \emph{2026-11-20:} The formal level restated: T14c deposited as version 1.0 on 20 November 2026 (10.5281/zenodo.24403592), after an independent review: stated at the particle model's density with its edges to R1 and R2; the bounds given as the ranges the deposited computation finds (P below 1 to 2e-16, e\^{}2/m above 7e7 to 9e7) and the surviving fraction set against T14's [...]
\item Theorem~\ref{c:dark-tower-matter.five-dimensional-towers} (\emph{The five-dimensional towers in cosmology and halos}), Comp 9.3 in the deposit: statement and proof. \emph{Store notes.} \emph{2026-11-19:} Leon's decision of 19 November 2026 (R3's unit is the length of Premise E): the statement used a depth of about a metre, this world's baryon spacing, on R3 alone. Assumes edges to R1 and R2 are added, since that value is the particle model's, and a marked note says so. No number changed. \emph{2026-11-20:} The formal level restated: T14c deposited as version 1.0 on 20 November 2026 (10.5281/zenodo.24403592), after an independent review: the dependence on the particle model's depth made a stated condition with edges to R1, R2 and R3 (it was a dated note); the plasma checks named I3 and I4; renumbered Computation 9.3. The technical level restated: as the formal [...]
\item Theorem~\ref{c:dark-tower-matter.r1-has-no-mechanism-odd-towers} (\emph{Neither E nor the particle model has a mechanism in the odd towers}), Obs 9.4 in the deposit: proof. \emph{Store notes.} \emph{2026-11-07:} Pin review: dark-matter.mirror-share-one-how-fragile repinned 6 -> 7 after reviewing its C1 restatement ((i) now rests on the common-ceiling-coupling premise, with an assumes edge, where it said P2-protected; the values are unchanged); the use here holds as cited. \emph{2026-11-20:} Title restated: T14c deposited as version 1.0 on 20 November 2026 (10.5281/zenodo.24403592), after an independent review: the result reaches premise E as well as the particle model, so the title names both; renumbered Observation 9.4. The formal level restated: as the title: nothing in the content fixes the masses or densities E equates; the bound on m/e\^{}2 [...]
\item Theorem~\ref{c:dark-tower-matter.population} (\emph{How the odd towers could be populated}), Obs 9.5 in the deposit: proof. \emph{Store notes.} \emph{2026-11-20:} The formal level restated: T14c deposited as version 1.0 on 20 November 2026 (10.5281/zenodo.24403592), after an independent review: the density to be filled is E's (was R2's); renumbered Observation 9.5. The technical level restated: as the formal level. The lay level restated: as the formal level. The one-liner level restated: as the formal level.
\item Proposition~\ref{c:dark-tower-matter.sterile-kinematics} (\emph{The kinematics of the sterile towers}), Prop 10.1 in the deposit: proof. \emph{Store notes.} \emph{2026-10-21:} Elementary signature facts proved directly and cited by nothing later (T14c is the last paper to treat the sterile towers); linked by lineage to the towers fact that identifies the sterile representatives as (2,2) and (4,0). \emph{2026-10-26:} Grounding pass 2026-10-26: cites synthesis.completion-recovers-the-real-clifford-algebras where the proof uses them.
\item Conjecture~\ref{c:dark-tower-matter.towers-built} (\emph{The dark towers as built, against the data}), Obs 11.1 in the deposit: statement. \emph{Store notes.} \emph{2026-10-21:} Terminal observation of T14c against the data, with full lineage to the results its table reads; nothing later cites it. \emph{2026-11-07:} G1 (DECISIONS; R05 item 2): the cold-mirror exception now says that it is an exception to P2 at its version-1 reach, that P2 as scoped (version 2, labels not states) leaves xi free (synthesis.tower-charge-sector-label-temperature-not), and that the cold mirror then needs a non-equilibrium descent (synthesis.cold-mirror-needs-non-equilibrium-descent); P2 is [...]
\end{itemize}
\subsection*{Reworded}
Wording, notation, terminology or citations only; the mathematics and the claim are as deposited.
\begin{itemize}\setlength{\itemsep}{2pt}
\item Remark~\ref{c:note.dark-tower-matter.branches-break-u-what-they-give} (\emph{Branches that break U, and what they give}), Rem 5.2 in the deposit: statement.
\item Proposition~\ref{c:dark-tower-matter.residue-seven-unit} (\emph{The central unit of the residue-seven tower}), Prop 7.1 in the deposit: proof.
\end{itemize}
\noindent In addition, 38 results from elsewhere in the series or the algebra core are stated here that the deposited paper cited or used without stating; 2 results gained citations only. These are not marked.
\begin{thebibliography}{99}
\bibitem{ButlerT14} L.~Butler,
\emph{Dark Matter and the Arrow of Time: The Five Dark Towers},
Transport Papers T14 (2026).
\bibitem{ButlerT9} L.~Butler,
\emph{The Twistor Correspondence: Ascended Spinors, Incidence, and the Palatial Bridge},
Transport Papers T9 (2026).
\bibitem{ButlerT3} L.~Butler,
\emph{The Tower and the CAR Limit: Iterated Transport, the Frozen Clock, and Second Quantisation},
Transport Papers T3 (2026), doi:\href{https://doi.org/10.5281/zenodo.22721556}{10.5281/zenodo.22721556}.
\bibitem{Chevalley1954} C.~Chevalley,
\emph{The Algebraic Theory of Spinors},
Columbia University Press (1954).
\bibitem{Lounesto2001} P.~Lounesto,
\emph{Clifford Algebras and Spinors}, 2nd ed.,
London Mathematical Society Lecture Note Series 286, Cambridge University Press (2001).
\bibitem{LawsonMichelsohn1989} H.~B. Lawson and M.-L. Michelsohn,
\emph{Spin Geometry},
Princeton University Press (1989).
\bibitem{ButlerT11} L.~Butler,
\emph{The Origin of Mass: The Weight Dictionary and the Exact Join},
Transport Papers T11 (2026), doi:\href{https://doi.org/10.5281/zenodo.22722565}{10.5281/zenodo.22722565}.
\bibitem{ButlerT12} L.~Butler,
\emph{The Standard Model from the Minimal Massive Composite},
Transport Papers T12 (2026), doi:\href{https://doi.org/10.5281/zenodo.22722669}{10.5281/zenodo.22722669}.
\bibitem{ButlerT14b} L.~Butler,
\emph{Simulating the Dark Towers: Methods, Assumptions and Tests of the Unit-Ball Composition},
Transport Papers T14b (2026), doi:\href{https://doi.org/10.5281/zenodo.23234576}{10.5281/zenodo.23234576}, deposited with its simulation code, run outputs and films; the films are also on YouTube, ``Galaxy clusters in collision'' (\url{https://www.youtube.com/watch?v=0SvmSwInN0c}) and ``The Bullet moment'' (\url{https://www.youtube.com/watch?v=KQnZUVR7bCk}).
\bibitem{Redlich1984} A.~N. Redlich,
\emph{Gauge noninvariance and parity nonconservation of three-dimensional fermions},
Phys.\ Rev.\ Lett.\ 52 (1984) 18--21; \emph{Parity violation and gauge noninvariance of the effective gauge field action in three dimensions}, Phys.\ Rev.\ D 29 (1984) 2366--2374.
\bibitem{IMS1997} K.~Intriligator, D.~R. Morrison, N.~Seiberg,
\emph{Five-dimensional supersymmetric gauge theories and degenerations of Calabi-Yau spaces},
Nucl.\ Phys.\ B 497 (1997) 56--100.
\bibitem{Widmark2021} A.~Widmark, C.~F.~P.~Laporte, P.~F.~de Salas, and G.~Monari,
\emph{Weighing the Galactic disk using phase-space spirals. II. Most stringent constraints on a thin dark disk using Gaia EDR3},
Astron. Astrophys. \textbf{653} (2021), A86.
\bibitem{RouxCline2020} J.-S.~Roux and J.~M.~Cline,
\emph{Constraining galactic structures of mirror dark matter},
Phys. Rev. D \textbf{102} (2020), 063518; arXiv:2001.11504.
\bibitem{Polyakov1977} A.~M. Polyakov,
\emph{Quark confinement and topology of gauge theories},
Nucl. Phys. B \textbf{120} (1977), 429--458.
\bibitem{DJT1982} \emph{Reference DJT1982 is not in the store's registry.}
\bibitem{AHPS1989} \emph{Reference AHPS1989 is not in the store's registry.}
\bibitem{Wilczek1982} \emph{Reference Wilczek1982 is not in the store's registry.}
\bibitem{CLP2000} \emph{Reference CLP2000 is not in the store's registry.}
\bibitem{SSS2024} \emph{Reference SSS2024 is not in the store's registry.}
\bibitem{PIC2024} \emph{Reference PIC2024 is not in the store's registry.}
\bibitem{Lasenby2020} \emph{Reference Lasenby2020 is not in the store's registry.}
\bibitem{ButlerT16} L.~Butler,
\emph{Quantum Mechanics and General Relativity from One Relation: The Two Shadows and the Laboratory},
Transport Papers T16 (2026), doi:\href{https://doi.org/10.5281/zenodo.22723212}{10.5281/zenodo.22723212}.
\bibitem{Planck2018} N.~Aghanim et al. (Planck Collaboration),
\emph{Planck 2018 results. VI. Cosmological parameters},
Astron. Astrophys. \textbf{641} (2020), A6.
\bibitem{ButlerT13} L.~Butler,
\emph{Quark and Lepton Families: Condensation, Democracy, and the Koide Relation},
Transport Papers T13 (2026), doi:\href{https://doi.org/10.5281/zenodo.22722845}{10.5281/zenodo.22722845}.
\bibitem{ButlerT15} L.~Butler,
\emph{Neutrino Masses: The Empty Twistor, the Seesaw Scale, and the Bet at the Floor},
Transport Papers T15 (2026), doi:\href{https://doi.org/10.5281/zenodo.22723070}{10.5281/zenodo.22723070}.
\bibitem{ButlerT14a} L.~Butler,
\emph{One Dimension at a Time: The Odd Extensions of the Oriented Towers},
Transport Papers T14a (2026), doi:\href{https://doi.org/10.5281/zenodo.23232104}{10.5281/zenodo.23232104}.
\bibitem{ButlerT17} L.~Butler,
\emph{The Transport Programme: Method, Status, and the Falsifiable Ledger},
Transport Papers T17 (2026).
\bibitem{ButlerVerifyT14c} \emph{Reference ButlerVerifyT14c is not in the store's registry.}
\bibitem{Porteous1995} I.~R.~Porteous,
\emph{Clifford Algebras and the Classical Groups},
Cambridge University Press (1995).
\bibitem{Hestenes1966} \emph{Reference Hestenes1966 is not in the store's registry.}
\bibitem{DoranLasenby2003} C.~Doran and A.~Lasenby,
\emph{Geometric Algebra for Physicists},
Cambridge University Press (2003).
\bibitem{FootMirror2004} R.~Foot,
\emph{Mirror matter-type dark matter},
Int. J. Mod. Phys. D \textbf{13} (2004), 2161--2192.
\end{thebibliography}
\end{document}